\documentclass[letterpaper]{article}
\usepackage{aaai}
\usepackage{times}
\usepackage{helvet}
\usepackage{courier}
\setlength{\pdfpagewidth}{8.5in}
\setlength{\pdfpageheight}{11in}
%\usepackage{amsthm}
\usepackage{amssymb}
\usepackage[utf8]{inputenc} % allow utf-8 input
\usepackage[T1]{fontenc}    % use 8-bit T1 fonts
\usepackage{hyperref}       % hyperlinks
\usepackage{url}            % simple URL typesetting
\usepackage{booktabs}       % professional-quality tables
\usepackage{amsfonts}       % blackboard math symbols
\usepackage{nicefrac}       % compact symbols for 1/2, etc.
\usepackage{microtype}      % microtypography
\usepackage{graphicx}
\usepackage{natbib}
\usepackage{doi}
\usepackage{amsmath}
\usepackage{mathtools}
\usepackage{physics}
\usepackage{wrapfig}
\usepackage{bm}
\usepackage{algorithm,algpseudocode}
\usepackage{xcolor}
\usepackage{forest}
\usepackage{tcolorbox}
\usepackage{tikz-cd}
\usepackage{enumitem}

% Keep section numbering on so Section~\ref{...} resolves to numbers.
\setcounter{secnumdepth}{2}

% Suppress the AAAI copyright slug for drafts.
\nocopyright


% A proof environment that accepts an optional argument (e.g. "Proof sketch")
% so that material from the merged hnn sources compiles without modification.
\makeatletter
\newenvironment{proof}[1][Proof]{\paragraph{#1:}}{\hfill$\square$}
\makeatother
\newenvironment{solution}{\paragraph{Solution:}}{\hfill$\square$}


%—– Blackboard‐bold sets —–
\newcommand{\R}{\mathbb{R}}
\newcommand{\N}{\mathbb{N}}
\newcommand{\Z}{\mathbb{Z}}
\newcommand{\Q}{\mathbb{Q}}
\newcommand{\C}{\mathbb{C}}
\newcommand{\D}{\mathbb{D}}      % Poincar\'e ball

%—– Probability & expectation —–
\newcommand{\Prb}{\mathbb{P}}        % probability
\newcommand{\E}{\mathbb{E}}           % expectation
\newcommand{\Var}{\mathrm{Var}}       % variance
\newcommand{\Cov}{\mathrm{Cov}}       % covariance

%—– Common math operators —–
\newcommand{\diag}{\operatorname{diag}}
\newcommand{\supp}{\operatorname{supp}}
\newcommand{\argmax}{\operatorname*{arg\,max}}
\newcommand{\argmin}{\operatorname*{arg\,min}}
\newcommand{\ind}{\mathbf{1}}         % indicator
% \tr is already provided by the physics package as \tr; do NOT redefine.

%—– Norms, inner products, absolute value —–
% \norm is provided by the physics package; we keep that.
% \inner is NOT provided by physics, so we define our own (matches the
% original template and is compatible with the merged hnn sources).
\newcommand{\inner}[2]{\left\langle #1, #2 \right\rangle}

%—– M\"obius / hyperbolic shorthands —–
\newcommand{\Mob}{\oplus_c}
\newcommand{\Mobm}{\otimes_c}
\newcommand{\exps}{\exp_{\mathbf{0}}^c}
\newcommand{\logs}{\log_{\mathbf{0}}^c}

%—– Calligraphic sets —–
\newcommand{\cA}{\mathcal{A}}
\newcommand{\cS}{\mathcal{S}}
\newcommand{\cT}{\mathcal{T}}
\newcommand{\cF}{\mathcal{F}}
\newcommand{\cE}{\mathcal{E}}

%—– RL / MDP notation —–
\newcommand{\SSS}{\mathcal{S}}        % state space
\newcommand{\AAA}{\mathcal{A}}        % action space
\newcommand{\PP}{\mathcal{P}}         % transition kernel
\newcommand{\RRR}{\mathcal{R}}        % reward function
\newcommand{\gammaRL}{\gamma}         % discount factor

%—– Confidence bounds —–
\newcommand{\olQ}{\overline{Q}}
\newcommand{\ulQ}{\underline{Q}}
\newcommand{\olV}{\overline{V}}
\newcommand{\ulV}{\underline{V}}

%—– Miscellaneous —–
\newcommand{\eps}{\varepsilon}
\newcommand{\dt}{\mathrm{d}t}
\newcommand{\given}{\,\bigl\vert\,}    % conditioning bar

%—– Shorthand for frequently used fractions —–
\newcommand{\half}{\tfrac12}
\newcommand{\thalf}{\tfrac32}

%—– Styling / margin notes —–
\newcommand{\todo}[1]{\textcolor{red}{[TODO: #1]}}
\newcommand{\william}[1]{{\color{purple}{[william: #1]}}}


\newcommand{\note}[1]{\begin{tcolorbox}
\textbf{Note:} #1\end{tcolorbox}}

\newtheorem{theorem}{Theorem}
\newtheorem{proposition}[theorem]{Proposition}
\newtheorem{lemma}[theorem]{Lemma}
\newtheorem{corollary}[theorem]{Corollary}
\newtheorem{definition}[theorem]{Definition}
\newtheorem{assumption}[theorem]{Assumption}
\newtheorem{remark}[theorem]{Remark}
\newtheorem{example}[theorem]{Example}
\newtheorem{conjecture}[theorem]{Conjecture}

% --- Allow long theorem/lemma/corollary notes to wrap ---------------------
% The article-class theorem head places "Name N (note)" inside a \item label,
% which is set in an unbreakable hbox and therefore overflows the (narrow,
% two-column) measure when the note is long. We redefine the heads so the
% note flows as ordinary breakable text.
\makeatletter
\renewcommand{\@begintheorem}[2]{%
  \trivlist
  \item\relax{\bfseries #1\ #2.}\hskip\labelsep\itshape}
\renewcommand{\@opargbegintheorem}[3]{%
  \trivlist
  \item\relax{\bfseries #1\ #2\ (#3).}\hskip\labelsep\itshape}
\makeatother
% --------------------------------------------------------------------------

\title{A title}
\author{William Chang\\
Department of Applied Mathematics\\
University of California, Los Angeles\\
Los Angeles, CA\\
\url{https://williamc.me/}\\
}

%%% Add PDF metadata to help others organize their library
%%% Once the PDF is generated, you can check the metadata with
%%% $ pdfinfo template.pdf
\hypersetup{
pdftitle={},
pdfsubject={},
pdfauthor={William Chang},
pdfkeywords={},
}

% =====================================================================
% BODY of restructured manuscript (v3)
%
% Two main results, both placed up front:
%   1. Curvature-adaptive transductive generalization bound for Q-GCN
%      (covering-number argument on a diffeomorphic product manifold).
%   2. Curvature-dependent single-step descent lemma for general
%      layerwise Riemannian neural networks, with recursively defined
%      smoothness constants L_net = H_N J_N^2 + G_N K_N, plus the
%      deep-linear specialization (Arora-style preconditioner with
%      curvature-dependent step size).
%
% Hyperbolic / pseudo-Riemannian specializations under tangent and
% intrinsic losses are now mentioned as future directions, not as
% standalone main results.
% =====================================================================

\begin{document}
\maketitle
\begin{abstract}
We study optimization and generalization for neural networks operating in
non-Euclidean geometry. Both of our main results live in a single model:
a \emph{general layerwise Riemannian neural network} whose $j$-th layer
maps $h_{j-1}\in\mathcal M_{j-1}$ to
$h_j = \exp_{o_j}^{(j)}(F_j(\log_{o_{j-1}}^{(j-1)}(h_{j-1});\theta_j))
\in\mathcal M_j$ on a complete Riemannian manifold $\mathcal M_j$, with
layerwise geometric constants
$\alpha_j,\beta_j,\gamma_j,\delta_j,P_j,Q_j$ derived from the Jacobians
and Hessians of the exponential and logarithmic maps and the tangent-space
operation $F_j$. We prove: \textbf{(i)} a \emph{curvature-dependent
single-step descent lemma} in which the smoothness constant
$L_{\mathrm{net}} = H_N J_N^2 + G_N K_N$ is built from recursively
defined layerwise quantities $J_j, K_j$ that propagate curvature bounds
through the network, yielding a curvature-modulated step-size restriction
$\eta\le 1/L_{\mathrm{net}}$. The descent lemma admits a deep-linear
specialization under approximate balancedness, recovering the
Arora--Cohen--Golowich--Hu \cite{arora2018convergence} single-step
inequality with a curvature-dependent step size. \textbf{(ii)} A
\emph{curvature-adaptive generalization bound} for the same model. The
bound depends on the input-manifold sectional curvature
$\kappa_{\max}^{(0)}$ through a Bishop--Gromov penalty
$\psi(\kappa_{\max}^{(0)}, L_\ell\Lambda_N) = \sqrt{|\kappa_{\max}^{(0)}|}/(L_\ell\Lambda_N)$,
and on every \emph{hidden}-manifold curvature through the network's
spatial Lipschitz constant
$\Lambda_N = \prod_{j=1}^N \alpha_j P_j \beta_{j-1}$. This is a layerwise
refinement of the recent manifold Rademacher bound of
\cite{sarkar2025learning}: their bound is recovered by treating the
network as a black-box Lipschitz function, ours is sharper whenever the
layerwise constants are individually controlled. As a corollary we
recover a transductive Q-GCN bound with an additive layerwise curvature
penalty $\Psi_{\mathrm{curv}} = \sum_\ell c_\ell \sqrt{|\kappa^{(\ell)}_{\max}|}$,
obtained by applying the layerwise theorem to the diffeomorphic product
manifold $M_\psi = \prod_\ell(\mathbb S^t_{-\beta_\ell}\times\R^{s_\ell})$
of \cite{xiong2022pseudo} in the transductive setting of
\cite{tang2023towards}. We close with a discussion of how the descent
lemma specializes to the tangent-space surrogate (where it recovers the
Euclidean rate verbatim), the intrinsic squared hyperbolic distance
(where curvature enters through $S_c(R) = \sinh(\sqrt c R)/(\sqrt c R)$),
and the pseudo-Riemannian setting under a spacelike-trajectory hypothesis.
\end{abstract}

% =====================================================================
\section{Introduction}
\label{sec:intro}
% =====================================================================

Hyperbolic and pseudo-Riemannian geometries have emerged as natural
ambient spaces for representations of tree-like, hierarchical, and
mixed-topology data. The exponential volume growth of hyperbolic space
admits low-distortion embeddings of trees in dimensions where Euclidean
embeddings provably cannot \cite{sala2018representation},
and this observation drove an influential line of work on hyperbolic
representation learning, from the Poincar\'e and Lorentz embeddings of
\cite{nickel2017poincare, nickel2018learning} to mixed-curvature product
spaces \cite{gu2018learning} that interpolate between hyperbolic,
spherical, and Euclidean factors. These representations have proven
effective across domains including word and language embeddings
\cite{tifrea2018poincar}, variational generative modeling
\cite{mathieu2019continuous, skopek2019mixed}, and computer vision
\cite{khrulkov2020hyperbolic, atigh2022hyperbolic}. Building on this,
\cite{tifrea2018poincar} introduced \emph{hyperbolic neural networks} ---
feed-forward and recurrent architectures operating natively on the
Poincar\'e ball --- and pseudo-Riemannian models such as the $Q$-GCN of
\cite{xiong2022pseudo} interpolate between hyperbolic and spherical
regimes within a single indefinite-signature manifold. Both lines have
seen substantial empirical uptake, but their theoretical behavior --- both
generalization and optimization --- is essentially unstudied.

The gap is striking given the maturity of the empirical literature.
Curvature-aware graph networks now span hyperbolic
\cite{yang2022hyperbolic}, constant-curvature
\cite{bachmann2020constant}, Lorentzian \cite{zhang2021lorentzian},
mixed-curvature \cite{sun2022self, wang2021mixed}, and
pseudo-Riemannian \cite{xiong2022pseudo, law2021ultrahyperbolic}
constructions, yet rigorous guarantees lag far behind. The handful of
theoretical results that do exist are narrow: \cite{suzuki2021generalization}
compare linear and hyperbolic graph-embedding generalization under
negative sampling, and \cite{li2024improving} use Lipschitz analysis to
study the robustness of hyperbolic networks, but neither addresses the
curvature dependence of either the generalization gap or the optimization
dynamics of a deep manifold-valued network. Two basic questions remain
open:

\begin{enumerate}[label=(\roman*), leftmargin=2.2em]
\item \emph{What is the generalization gap of a non-Euclidean graph neural
network, as a function of its curvature?} The Euclidean baseline is the
transductive generalization bound of \cite{tang2023towards}; its proof
depends on the covering geometry of a Euclidean function class. For a graph
network that operates through diffeomorphic charts onto a curved product
manifold, the covering argument must be replaced by a curvature-adaptive
one.

\item \emph{Does gradient-based optimization of a deep neural network
operating on a curved manifold admit a curvature-aware descent guarantee,
and how does layerwise curvature propagate into the admissible step
size?} The Euclidean precedent is the linear-rate convergence theorem of
\cite{arora2018convergence} for deep linear networks, whose proof rests on
the quadratic structure of the loss as a function of the end-to-end matrix
and on a deficiency-margin condition. In a manifold-valued network, the
chain rule introduces Jacobians of the layerwise exponential and
logarithmic maps, and the resulting smoothness constant must absorb all of
the layerwise curvatures in a controlled way.
\end{enumerate}

This paper addresses both questions for the architectures named above.
Section~\ref{sec:relwork} reviews the relevant prior work,
Section~\ref{sec:contrib} states our contributions,
Section~\ref{sec:prelim} fixes notation and develops the necessary
geometric machinery (Q-GCN architecture, general layerwise Riemannian
neural network, recursive smoothness constants), and
Section~\ref{sec:main} states the main results, organized by model
class (linear case, nonlinear case, graph case), with the discussion of
each result given in place and all proofs deferred to
Appendix~\ref{sec:proofs}. We close with future directions in
Section~\ref{sec:future}.

% =====================================================================
\section{Related Work}
\label{sec:relwork}
% =====================================================================

\paragraph{Non-Euclidean representation learning.}
The case for curved geometry begins with embeddings. Trees and other
hierarchical metrics embed into the hyperbolic plane with arbitrarily
small distortion \cite{sarkar2025learning}, and the precision-versus-dimension
tradeoffs of such embeddings were quantified by
\cite{sala2018representation}. \cite{nickel2017poincare, nickel2018learning}
showed that learning Poincar\'e and Lorentz embeddings recovers latent
hierarchies in symbolic and network data, while \cite{gu2018learning}
extended the idea to products of constant-curvature factors so that a
single space can match mixed graph topologies. Subsequent work carried
hyperbolic representations into language \cite{tifrea2018poincar},
generative modeling \cite{mathieu2019continuous, skopek2019mixed}, and
vision \cite{khrulkov2020hyperbolic, atigh2022hyperbolic}. A recurring
practical obstacle is numerical: the unbounded metric near the boundary
of the Poincar\'e ball makes embeddings ill-conditioned, motivating
tiling-based \cite{yu2019numerically} and stability-focused
\cite{mishne2023numerical} reformulations. These difficulties are part of
why a curvature-explicit theory of the downstream networks --- the subject
of this paper --- is valuable.

\paragraph{Hyperbolic and pseudo-Riemannian neural architectures.}
\cite{tifrea2018poincar} introduced M\"obius matrix multiplication and
exponential/logarithmic maps for the Poincar\'e ball, with empirical
validation on entailment and language tasks; this is one of the two
instantiations of the general layerwise Riemannian network we analyze.
Fully hyperbolic formulations that operate intrinsically in the Lorentz
model \cite{chen2022fully} and their convolutional
\cite{bdeir2024fully}, residual \cite{van2023poincare}, and attention
\cite{gulcehre2018hyperbolic} variants followed.
\cite{law2021ultrahyperbolic} moved beyond definite curvature to
ultrahyperbolic (pseudo-Riemannian) networks, and \cite{xiong2022pseudo}
introduced the Q-GCN architecture; the spherical projection $\psi$ and the
diffeomorphic $\log/\exp$ construction used in our generalization bound
come from that paper. The robustness of these architectures was recently
studied through Lipschitz analysis by \cite{li2024improving}, whose
layerwise Lipschitz viewpoint is complementary to the layerwise Lipschitz
constants $\alpha_j, \beta_j, P_j$ that drive both of our main results.

\paragraph{Graph neural networks in non-Euclidean geometry.}
Euclidean message passing \cite{kipf2016semi, hamilton2017inductive,
velivckovic2017graph, xu2018powerful} was lifted to curved spaces by a
large body of work. Hyperbolic graph convolutions
\cite{yang2022hyperbolic}, their constant-curvature
\cite{bachmann2020constant}, Lorentzian \cite{zhang2021lorentzian}, and
hyperbolic-to-hyperbolic \cite{dai2021hyperbolic} refinements, dynamic
\cite{sun2021hyperbolic} and adaptive-curvature \cite{fu2021ace}
variants, and the mixed-curvature \cite{sun2022self, wang2021mixed,
cho2023curve}, spectral \cite{grover2025spectro}, and spiking
\cite{sun2024spiking} extensions are surveyed in
\cite{yang2022hyperbolic}. The pseudo-Riemannian Q-GCN of
\cite{xiong2022pseudo} is the model to which our transductive
generalization corollary
(Corollary~\ref{thm:prgcn_transductive_generalization}) specializes.

\paragraph{Optimization of deep linear networks.}
Trajectory-based analysis of gradient descent for deep linear networks was
initiated by \cite{saxe2013exact} and developed through a sequence of
landscape and convergence results: the absence of poor local minima
\cite{lu2017depth, lu2017depth, laurent2018deep}, the role of
identity/residual structure \cite{hardt2016identity}, and the implicit
acceleration induced by overparameterization
\cite{arora2018convergence}. \cite{bartlett2017spectrally} proved
linear-rate convergence under identity initialization for linear residual
networks, and \cite{arora2018convergence} extended this to general
configurations under approximate balancedness and a deficiency margin;
related convergence guarantees appear in
\cite{shamir2019exponential, du2019width, du2019gradient,
allen2019convergence, zou2020global, ji2018gradient, nouiehed2022learning},
and a Riemannian gradient-flow account is given by \cite{bah2022learning}.
The factorization-induced preconditioner of
\cite{arora2018convergence, arora2018convergence} is the algebraic core
of the deep-linear specialization of our descent lemma
(Theorem~\ref{thm:linear-descent}).

\paragraph{Riemannian optimization.}
A separate tradition optimizes \emph{over} a manifold-valued parameter:
the foundational treatments of \cite{absil2008optimization} and
\cite{boumal2023introduction} (see also the survey of
\cite{hu2020brief}), stochastic and variance-reduced methods
\cite{bonnabel2013stochastic, zhang2016riemannian, sato2019riemannian},
geodesically convex first-order analysis \cite{zhang2016first}, and the
Riemannian adaptive methods of \cite{bécigneul2019riemannianadaptiveoptimizationmethods} used to
train hyperbolic models. Our setting is dual to this one: we optimize a
\emph{Euclidean} parameter under a loss whose landscape is shaped by a
non-Euclidean \emph{target} geometry, so the curvature enters through the
Jacobians and Hessians of the layerwise $\exp$ and $\log$ maps rather than
through a Riemannian metric on the parameter space. Convergence under
Riemannian losses has otherwise been studied mainly in shallow matrix
factorization and completion
\cite{sun2016guaranteed}; to
our knowledge no prior work derives a single-step descent inequality for
\emph{deep} manifold-valued networks with curvature-dependent smoothness
constants.

\paragraph{Generalization of graph neural networks.}
Generalization theory for Euclidean GNNs has been approached through VC
dimension \cite{scarselli2018vapnik}, algorithmic stability
\cite{verma2019stability}, Rademacher complexity
\cite{garg2020generalization}, and PAC-Bayes
\cite{liao2020pac, ju2023generalization, sun2024pac}, with further
analyses for large random graphs \cite{maskey2022generalization} and
learning-theoretic explanations of empirical behavior
\cite{esser2021learning}; the transductive bound of \cite{tang2023towards}
is the Euclidean baseline whose covering geometry our Q-GCN corollary
replaces with a curvature-adaptive one. These build on the broader
norm-based and size-independent generalization theory for deep networks
\cite{bartlett2017spectrally, neyshabur2017pac, golowich2018size,
long2020generalizationboundsdeepconvolutional}. Hyperbolic generalization
has barely been touched: \cite{suzuki2021generalization} compare linear
and hyperbolic graph embeddings under negative sampling, but do not treat
the curvature dependence of a deep network's generalization gap, which is
the object of our Theorem~\ref{thm:layerwise-generalization}.

\paragraph{Generalization on Riemannian manifolds.}
Curvature-adaptive covering-number bounds on Riemannian manifolds go back
to \cite{niyogi2008finding} via Bishop--Gromov volume comparison, where
reach and curvature control the covering numbers used throughout
manifold-learning theory \cite{genovese2012minimax, kim2015tight,
fefferman2016testing}. \cite{sarkar2025learning} recently used these to
give a Rademacher generalization bound for a single $L$-Lipschitz function
on a compact Riemannian manifold, with an explicit Bishop--Gromov
curvature penalty $\psi(\kappa,L) = \sqrt{|\kappa|}/L$. In their setting
the network is treated as a single Lipschitz function, so the curvature of
any hidden representation does not enter the bound. Our
Theorem~\ref{thm:layerwise-generalization} keeps their curvature-adaptive
covering of the input manifold but \emph{layers} the Lipschitz analysis,
so that the curvature of every hidden manifold propagates into the bound
through the same recursive constants $\alpha_j, \beta_j, P_j$ used by our
descent lemma. The Euclidean GNN baseline of \cite{tang2023towards} is the
transductive setting we use to specialize to Q-GCN in
Corollary~\ref{thm:prgcn_transductive_generalization}.

% =====================================================================
\paragraph{Our Contributions}
\label{sec:contrib}
% =====================================================================

We make three contributions, organized by model class to match the
presentation in Section~\ref{sec:main}.

\begin{itemize}[leftmargin=1.4em, itemsep=4pt]
\item \textbf{Linear case} (Theorem~\ref{thm:linear-descent},
Section~\ref{sec:main-linear}). For a deep linear network on a Riemannian
output manifold --- every tangent-space map linear, $F_j(u) = W_j u$, so
the network collapses to a single end-to-end matrix product
$A = W_N \cdots W_1$ wrapped by $\exp_{o_N}^{(N)}$ --- we prove a
single-step descent inequality of Arora--\allowbreak Cohen--\allowbreak
Golowich--\allowbreak Hu form under a
curvature-dependent step-size restriction, with the output-manifold
curvature entering through the smoothness constant $L_{\kappa,R,D}$.

\item \textbf{Nonlinear case} (Theorems~\ref{thm:general-descent}
and~\ref{thm:layerwise-generalization}, Section~\ref{sec:main-nonlinear}).
For the general layerwise Riemannian network, whose $j$-th layer maps
$h_{j-1}\in\mathcal M_{j-1}$ to
$h_j = \exp_{o_j}^{(j)}(F_j(\log_{o_{j-1}}^{(j-1)}(h_{j-1}); \theta_j))$
with possibly nonlinear $F_j$, we prove (i) a single-step descent lemma
whose admissible step size $\eta \le 1/L_{\mathrm{net}}$ is controlled by
the recursively defined smoothness constant
$L_{\mathrm{net}} = H_N J_N^2 + G_N K_N$, and (ii) a curvature-adaptive
generalization bound whose network spatial Lipschitz constant
$\Lambda_N = \prod_{j=1}^N \alpha_j P_j \beta_{j-1}$ is expressed
layerwise through the same constants. Both are built from the layerwise
Jacobian/Hessian bounds $\alpha_j, \beta_j, \gamma_j, \delta_j, P_j$ of
the exponential and logarithmic maps. The generalization bound refines
\cite{sarkar2025learning}, who treat the network as a black-box Lipschitz
function on $\mathcal M_0$ and hence see only the input-manifold
curvature; ours additionally exposes every hidden-manifold curvature
through $\Lambda_N$.

\item \textbf{Graph special case}
(Corollary~\ref{thm:prgcn_transductive_generalization},
Section~\ref{sec:main-graph}). Applying
Theorem~\ref{thm:layerwise-generalization} to the diffeomorphic product
manifold $M_\psi$ of \cite{xiong2022pseudo}, with the standard PAC
concentration replaced by the transductive Rademacher comparison of
\cite{tang2023towards}, we obtain a curvature-adaptive transductive
generalization bound for Q-GCN with an additive layerwise curvature
penalty $\Psi_{\mathrm{curv}} = \sum_\ell c_\ell
\sqrt{|\kappa^{(\ell)}_{\max}|}$ arising from applying Bishop--Gromov
factorwise on $M_\psi$.
\end{itemize}

Each main result is followed by its own discussion in place. Future
directions are collected in Section~\ref{sec:future}.

% =====================================================================
\section{Preliminaries}
\label{sec:prelim}
% =====================================================================

This section fixes notation and isolates the geometric machinery used in
the main results. We first recall the Euclidean GNN baseline studied in
\cite{tang2023towards} and the pseudo-Riemannian Q-GCN of
\cite{xiong2022pseudo} (Section~\ref{sec:eucl-gnn}--\ref{sec:prelim-qgcn}).
We then introduce the general layerwise Riemannian neural network model
that underlies the descent lemma (Section~\ref{sec:prelim-riem-nn}),
develop the recursive smoothness constants $J_j, K_j, L_{\mathrm{net}}$
(Section~\ref{sec:prelim-recursive}), and finally state the two structural
conditions of \cite{arora2018convergence} that drive the deep-linear
specialization (Section~\ref{sec:prelim-tools}).

\subsection{Euclidean GNN baseline}
\label{sec:eucl-gnn}

Let $G = (\mathcal V, \mathcal E)$ be a graph with $n$ nodes, node-feature
matrix $X \in \R^{n \times d}$, adjacency $A$, degree matrix $D$, and
normalized adjacency $\tilde A = (D+I)^{-1/2}(A+I)(D+I)^{-1/2}$. A two-layer
GCN takes the form
$\hat Y = \mathrm{Softmax}(\tilde A \sigma(\tilde A X W_1) W_2)$,
where $W_1 \in \R^{d \times h}$, $W_2 \in \R^{h \times |\mathcal Y|}$, and
$\sigma$ is a pointwise nonlinearity. The Euclidean model alternates three
operations — linear transformation, graph aggregation, and pointwise
nonlinearity — and this decomposition determines the architecture-dependent
constants in the generalization bound of \cite{tang2023towards}. The
parameter-Lipschitz constant of the per-sample loss is controlled by
Jacobians of its constituent operations; this mechanism — \emph{parameter
Lipschitz constant inherited from Jacobians of the constituent
operations} — is what carries over to the pseudo-Riemannian case, with the
new terms coming from the Jacobians of $\widehat{\log}$, $\widehat{\exp}$,
and parallel transport.

\subsection{Pseudo-Riemannian Q-GCN architecture}
\label{sec:prelim-qgcn}

We now summarize the pseudo-Riemannian graph convolutional network of
\cite{xiong2022pseudo}. The essential difference from the Euclidean
baseline is that hidden states no longer evolve in a Euclidean vector
space but on a pseudo-hyperboloid, and aggregation is performed in a
tangent space.

\paragraph{Pseudo-hyperboloid.}
Let $\R^{s, t+1}$ denote $\R^{s+t+1}$ equipped with the indefinite bilinear
form $\langle x, y \rangle_t = -\sum_{i=0}^t x_i y_i + \sum_{j=t+1}^{s+t}
x_j y_j$. The \emph{pseudo-hyperboloid} is $\mathcal Q_\beta^{s,t} = \{x
\in \R^{s,t+1} : \langle x, x \rangle_t = \beta\}$ for $\beta < 0$. This
geometry contains hyperbolic and spherical models as special or induced
cases and is flexible enough to represent mixed graph topologies.

\paragraph{Diffeomorphic log/exp.}
$\mathcal Q_\beta^{s,t}$ is not globally geodesically connected: for some
$(x, y)$, the usual logarithmic map $\log_x(y)$ is undefined. Q-GCN
sidesteps this by transporting the geometry through a diffeomorphism
$\psi : \mathcal Q_\beta^{s,t} \to \mathbb S^t_{-\beta} \times \R^s$ to a
geodesically well-behaved product manifold. The resulting diffeomorphic
logarithmic and exponential maps are
\[
\begin{aligned}
\widehat{\log}^c_{\mathcal Q_\beta}(x) &:=
\psi^{-1}\!\big(\log_{\mathbb S^t_{-\beta} \times \R^s}(\psi(x))\big),\\
\widehat{\exp}^d_{\mathcal Q_\beta}(\xi) &:=
\psi^{-1}\!\big(\exp_{\mathbb S^t_{-\beta} \times \R^s}(\psi(\xi))\big).
\end{aligned}
\]

\paragraph{Lifted linear maps and bias translation.}
Fix a reference point $o \in \mathcal Q_\beta^{s,t}$. Euclidean maps are
lifted to the manifold by conjugation with $\widehat{\log}_o$ and
$\widehat{\exp}_o$:
$W_\ell \otimes_{\beta_\ell} h := \widehat{\exp}^d_{o, \beta_\ell}(W_\ell
\widehat{\log}^c_{o, \beta_\ell}(h))$. Bias addition uses parallel transport:
$\tilde h \oplus_{\beta_\ell} b_\ell := \exp^{\beta_\ell}_{\tilde h}(
P^{\beta_\ell}_{o \to \tilde h}(b_\ell))$, with a special antipodal branch
when the direct geodesic is broken.

\paragraph{Layer update.}
For node $i$ at layer $\ell$ with embedding $h_i^\ell \in \mathcal
Q_{\beta_\ell}^{s,t}$, set $\bar h_j^\ell = (W_\ell \otimes_{\beta_\ell}
h_j^\ell) \oplus_{\beta_\ell} b_\ell$. The Q-GCN update is
\[
h_i^{\ell+1} = \widehat{\exp}^d_{o, \beta_{\ell+1}}\!\bigg(
\sigma\!\bigg(\sum_{j \in \mathcal N(i) \cup \{i\}}
\widehat{\log}^c_{o, \beta_\ell}(\bar h_j^\ell)\bigg)\bigg).
\]
Euclidean features $x_i \in \R^{d_0}$ are mapped onto the
pseudo-hyperboloid via a smooth initialization $h_i^0 = \varphi(x_i)$; the
final embedding is mapped back to a tangent space and classified by a
Euclidean head $\mathrm{logits}_i = W_{\mathrm{cls}} \widehat{\log}^c_{o,
\beta_L}(h_i^L)$.

\subsection{General layerwise Riemannian neural network}
\label{sec:prelim-riem-nn}

We now define the model that underlies our descent lemma. The Q-GCN
architecture above is one instance of it; the linear hyperbolic networks of
\cite{tifrea2018poincar} are another.

Let \(N\geq 1\) be the number of layers. For each layer \(j=0,\dots,N\),
let \((\mathcal M_j,g_j)\) be a complete Riemannian manifold of dimension
\(d_j\), with base point \(o_j\in \mathcal M_j\). Let
\(V_j:=T_{o_j}\mathcal M_j\simeq \mathbb R^{d_j}\) after choosing an
orthonormal frame. The data are
\[
(x_i,y_i)_{i=1}^m,
\qquad
x_i\in \mathcal M_0,
\qquad
y_i\in \mathcal M_N.
\]
For each layer \(j=1,\dots,N\), let \(F_j:V_{j-1}\times \Theta_j\to V_j\)
be a smooth tangent-space operation with parameters
\(\theta_j\in \Theta_j\). A basic example is \(F_j(u;\theta_j) =
\sigma_j(W_j u+b_j)\), where \(W_j\in \mathbb R^{d_j\times d_{j-1}}\) and
\(b_j\in V_j\); thus \(W_j\) is a Euclidean linear map \(V_{j-1}\to V_j\),
not a direct map between manifolds.

The corresponding manifold layer is \(\Phi_j:\mathcal M_{j-1}\times
\Theta_j\to \mathcal M_j\) defined by
\begin{equation}
\label{eq:phi-layer}
\Phi_j(h;\theta_j)
:=
\exp_{o_j}^{(j)}
\left(
F_j\left(
\log_{o_{j-1}}^{(j-1)}(h);
\theta_j
\right)
\right).
\end{equation}
Therefore the forward pass is \(h_0^{(i)}:=x_i\), and for \(j=1,\dots,N\),
\(u_j^{(i)}:=\log_{o_{j-1}}^{(j-1)}(h_{j-1}^{(i)})\in V_{j-1}\),
\(a_j^{(i)}:=F_j(u_j^{(i)};\theta_j)\in V_j\),
\(h_j^{(i)}:=\exp_{o_j}^{(j)}(a_j^{(i)})\in \mathcal M_j\). Equivalently,
$h_N^{(i)}(\theta) = \Phi_N(\cdot;\theta_N) \circ \cdots \circ
\Phi_1(\cdot;\theta_1)(x_i)$ with $\Phi_j$ as
in~\eqref{eq:phi-layer}. Let \(\theta:=(\theta_1,\dots,\theta_N)\)
denote the full parameter vector. We train using the intrinsic output loss
\[
\begin{gathered}
\mathcal L(\theta)
:=
\frac1m\sum_{i=1}^m
\ell_i(h_N^{(i)}(\theta)),\\
\ell_i(p)
:=
\frac12 d_N(p,y_i)^2,
\quad p\in \mathcal M_N.
\end{gathered}
\]

\paragraph{Layerwise geometric tube.}
For each \(j=0,\dots,N\), fix a compact tube \(\mathcal T_j\subset
\mathcal M_j\) such that all hidden states remain inside it:
\(h_j^{(i)}(\theta)\in \mathcal T_j\) for all \(i,j\). Assume
\(\mathcal T_j\) is contained in a normal neighborhood of \(o_j\), so that
\(\log_{o_j}^{(j)}\) is smooth on \(\mathcal T_j\); also assume that
\(\mathcal T_N\) avoids the cut locus of every \(y_i\), so that
\(p\mapsto \frac12 d_N(p,y_i)^2\) is smooth on \(\mathcal T_N\).

\paragraph{Layerwise geometric constants.}
Define
\[
\begin{aligned}
\alpha_j
&:=
\sup_{\|v\|\le R_j}
\|D\exp_{o_j}^{(j)}(v)\|_{\mathrm{op}},\\
\beta_j
&:=
\sup_{h\in \mathcal T_j}
\|D\log_{o_j}^{(j)}(h)\|_{\mathrm{op}},\\
\gamma_j
&:=
\sup_{\|v\|\le R_j}
\|D^2\exp_{o_j}^{(j)}(v)\|_{\mathrm{op}},\\
\delta_j
&:=
\sup_{h\in \mathcal T_j}
\|D^2\log_{o_j}^{(j)}(h)\|_{\mathrm{op}}.
\end{aligned}
\]
where \(R_j\) is chosen so that \(a_j^{(i)}(\theta)\in B_{V_j}(0,R_j)\) for
all \(i,j\). If \(\mathcal M_j=\mathbb H^{d_j}_{-\kappa_j}\), one may take
\[
\begin{aligned}
\alpha_j
&\lesssim
S_{\kappa_j}(R_j),
\qquad\text{where}\\
S_{\kappa_j}(R_j)
&:=
\begin{cases}
\dfrac{\sinh(\sqrt{\kappa_j}R_j)}
{\sqrt{\kappa_j}R_j}, & \kappa_j>0,\\[1.2em]
1, & \kappa_j=0.
\end{cases}
\end{aligned}
\]
Thus the curvature of every hidden manifold \(\mathcal M_j\) enters through
the constants \(\alpha_j,\beta_j,\gamma_j,\delta_j\). For each \(F_j\),
define the tangent-map constants
\(P_j := \sup\|D_uF_j(u;\theta_j)\|_{\mathrm{op}}\),
\(Q_j := \sup\|D^2_{uu}F_j(u;\theta_j)\|_{\mathrm{op}}\),
\(R_j^{\theta} := \sup\|D_{\theta_j}F_j(u;\theta_j)\|_{\mathrm{op}}\),
\(S_j^{\theta} := \sup\|D^2_{u\theta_j}F_j(u;\theta_j)\|_{\mathrm{op}}\),
and \(T_j^{\theta} := \sup\|D^2_{\theta_j\theta_j}F_j(u;\theta_j)\|_{\mathrm{op}}\),
with suprema taken over the relevant tangent-space tubes and parameter
region.

\subsection{Recursive smoothness constants}
\label{sec:prelim-recursive}

We now define recursively two quantities \(J_j\) and \(K_j\), controlling
the first and second derivatives of the hidden state \(h_j^{(i)}(\theta)\)
with respect to the full parameter vector \(\theta\). These are the
\emph{recursive lemma} that drives the descent statement of
Theorem~\ref{thm:general-descent}.

Since \(h_0^{(i)}=x_i\) is independent of \(\theta\), set \(J_0:=0\),
\(K_0:=0\). Inductively, given \(J_{j-1}, K_{j-1}\): since
\(u_j^{(i)} = \log_{o_{j-1}}^{(j-1)}(h_{j-1}^{(i)})\), we have
\(\|D_\theta u_j^{(i)}\| \le \beta_{j-1}J_{j-1}\) and
\(\|D^2_\theta u_j^{(i)}\| \le \delta_{j-1}J_{j-1}^2+\beta_{j-1}K_{j-1}\).
Next, since \(a_j^{(i)}=F_j(u_j^{(i)};\theta_j)\), define
\[
\widetilde J_j
:=
P_j\beta_{j-1}J_{j-1}
+
R_j^\theta,
\]
and
\[
\begin{aligned}
\widetilde K_j
:=\;
& Q_j(\beta_{j-1}J_{j-1})^2
+
P_j(\delta_{j-1}J_{j-1}^2+\beta_{j-1}K_{j-1})\\
& +
2S_j^\theta\beta_{j-1}J_{j-1}
+
T_j^\theta,
\end{aligned}
\]
so that \(\|D_\theta a_j^{(i)}\| \le \widetilde J_j\) and
\(\|D^2_\theta a_j^{(i)}\| \le \widetilde K_j\).
Finally, since \(h_j^{(i)} = \exp_{o_j}^{(j)}(a_j^{(i)})\),
\begin{equation}
\label{eq:JK-recursive}
J_j := \alpha_j\widetilde J_j,
\qquad
K_j := \gamma_j\widetilde J_j^2 + \alpha_j\widetilde K_j.
\end{equation}
Then \(\|D_\theta h_j^{(i)}\|\le J_j\) and
\(\|D^2_\theta h_j^{(i)}\|\le K_j\).

\paragraph{Output loss constants.}
For \(\ell_i(p)=\frac12 d_N(p,y_i)^2\), define
$G_N := \sup_{i,p\in\mathcal T_N}\|\nabla_p \ell_i(p)\|$ and
$H_N := \sup_{i,p\in\mathcal T_N}\|\nabla_p^2 \ell_i(p)\|_{\mathrm{op}}$.
If \(d_N(p,y_i)\le D_N\) on \(\mathcal T_N\), then \(G_N\le D_N\). For
hyperbolic output curvature \(-\kappa_N\),
\(H_N \lesssim \max\{1,\, D_N\sqrt{\kappa_N}\coth(\sqrt{\kappa_N}D_N)\}\).

\paragraph{Network smoothness constant.}
Define
\begin{equation}
\label{eq:Lnet-def}
L_{\mathrm{net}}
:=
H_NJ_N^2
+
G_NK_N.
\end{equation}
This constant depends on all layerwise curvature quantities
\(\{\alpha_j,\beta_j,\gamma_j,\delta_j\}_{j=0}^N\), and therefore on all
curvatures \(\kappa_0,\kappa_1,\dots,\kappa_N\) in the constant-curvature
case.

\subsection{Background tools: balancedness and deficiency margin}
\label{sec:prelim-tools}

The two structural conditions of \cite{arora2018convergence} appear in the
deep-linear specialization of our descent lemma.

\begin{definition}[$\delta$-balancedness]
\label{def:balanced}
Weight matrices $W_1, \ldots, W_L$ are \emph{$\delta$-balanced} if
$\|W_{\ell+1}^\top W_{\ell+1} - W_\ell W_\ell^\top\|_F \leq \delta$ for all
$1 \leq \ell \leq L-1$.
\end{definition}

\begin{definition}[Deficiency margin]
\label{def:defmargin}
A matrix $W \in \R^{d_L \times d_0}$ has \emph{deficiency margin $c' > 0$}
with respect to $\Phi$ if $\|W - \Phi\|_F \leq \sigma_{\min}(\Phi) - c'$.
We write $c'$ to avoid collision with the curvature parameter $c$.
\end{definition}

Approximate balancedness is automatically met by residual/identity
initialization and, with high probability, by suitably scaled near-zero
initializations \cite{arora2018convergence}. Deficiency margin ensures the
initial end-to-end matrix is not near a rank-deficient solution. Both
conditions are necessary for linear convergence in the Euclidean baseline
\cite[Appendix~C]{arora2018convergence}.


% =====================================================================
\section{Main Results}
\label{sec:main}
% =====================================================================

We organize the results by model class. Section~\ref{sec:main-linear}
treats the \emph{linear case}: a deep linear network on a Riemannian
output manifold, for which we obtain a single-step descent inequality of
Arora--Cohen--Golowich--Hu form with a curvature-dependent step size.
Section~\ref{sec:main-nonlinear} treats the \emph{nonlinear case}: a
general layerwise Riemannian neural network with possibly nonlinear
tangent-space operations, for which we obtain both a single-step descent
lemma and a curvature-adaptive generalization bound built from the same
recursive layerwise constants.
Section~\ref{sec:main-graph} treats the \emph{graph} setting as a
special case: the pseudo-Riemannian Q-GCN of \cite{xiong2022pseudo}
appears as a corollary of the nonlinear-case generalization bound.
All proofs are deferred to Appendix~\ref{sec:proofs} and the reader is
referred there as the results are stated.

% =====================================================================
\subsection{Linear case}
\label{sec:main-linear}
% =====================================================================

We begin with the simplest case: a deep linear network on a Riemannian
output manifold. Every tangent-space map is purely linear,
\(F_j(u;\theta_j)=W_j u\), with no biases, activations, skip
connections, or aggregations. Then
\(\log_{o_j}^{(j)}(h_j) = W_j \log_{o_{j-1}}^{(j-1)}(h_{j-1})\), so
the network collapses to
\[
h_N
=
\exp_{o_N}^{(N)}\!\left(
W_N\cdots W_1\log_{o_0}^{(0)}(x)
\right).
\]
If every tangent map is purely linear, \(F_j(u;\theta_j)=W_ju\), and there
are no biases, activations, skip connections, or aggregations, then
\(\log_{o_j}^{(j)}(h_j) = W_j\log_{o_{j-1}}^{(j-1)}(h_{j-1})\), so the
network collapses to
\[
h_N
=
\exp_{o_N}^{(N)}
\left(
W_N\cdots W_1\log_{o_0}^{(0)}(x)
\right).
\]
We introduce the Arora-style end-to-end matrix
\(A=W_N\cdots W_1\) and the balancedness condition
\(\|W_{j+1}^\top W_{j+1}-W_jW_j^\top\|_F\le \delta\) of
Definition~\ref{def:balanced}, and obtain a curvature-aware analogue of the
Arora--Cohen--Golowich--Hu \cite{arora2018convergence} single-step
inequality: a per-step contraction of the end-to-end loss by the
preconditioned factor
$\sigma_{\min}(A)^{2(N-1)/N}\|\nabla_A \mathcal F(A)\|_{\mathrm{HS}}^2$,
which yields linear-rate convergence under iteration. The precise
statement is as follows.

For \(1\leq a\leq b\leq N\), write \(W_{a:b}:=W_bW_{b-1}\cdots W_a\), with
the convention that empty products are identity maps.

\begin{theorem}[Riemannian single-step descent for a deep linear network]
\label{thm:linear-descent}
Let \(A=W_{1:N}=W_N\cdots W_1\). Assume that \(A\in\mathcal T_{R,D}\), and
let \(G:=\nabla_A \mathcal F(A)\). Suppose:
\begin{enumerate}[leftmargin=2em]
\item The layers are uniformly bounded:
\(\|W_j\|_{\mathrm{op}}\leq M\), \(j=1,\dots,N\).
\item The end-to-end map is nondegenerate: \(\sigma_{\min}(A)>0\).
\item The factors are sufficiently balanced so that the factor-gradient
coercivity bound holds:
\[
\begin{aligned}
&\sum_{j=1}^N
\|W_{j+1:N}^*GW_{1:j-1}^*\|_{\mathrm{HS}}^2\\
&\qquad\geq
\tfrac34 \sigma_{\min}(A)^{2(N-1)/N} \|G\|_{\mathrm{HS}}^2.
\end{aligned}
\]
\item The update is small enough that the segment between \(A\) and
\(A^+\) remains in \(\mathcal T_{R,D}\).
\end{enumerate}
Let gradient descent on the Euclidean/tangent-space parameters be
\(W_j^+ = W_j - \eta \nabla_{W_j}\mathcal L(W_1,\dots,W_N)\), and let
\(A^+:=W_N^+\cdots W_1^+\). There exists a constant \(C_N>0\), depending
only on depth \(N\), such that if
\begin{equation}
\label{eq:linear-stepsize}
\begin{aligned}
\eta C_N &\left( M^{3N-4}G_{\kappa,R,D} + L_{\kappa,R,D}M^{4N-4} \right)\\
&\leq \tfrac14 \sigma_{\min}(A)^{2(N-1)/N},
\end{aligned}
\end{equation}
where \(G_{\kappa,R,D} := \sup_{B\in\mathcal T_{R,D}}\|\nabla_A\mathcal F(B)\|_{\mathrm{HS}}\)
and \(L_{\kappa, R, D}\) is the curvature-dependent smoothness constant of
$\mathcal F$ on the tube, then
\begin{equation}
\label{eq:linear-descent}
\mathcal F(A^+)-\mathcal F(A)
\leq
- \tfrac{\eta}{2} \sigma_{\min}(A)^{2(N-1)/N}
\|\nabla_A\mathcal F(A)\|_{\mathrm{HS}}^2.
\end{equation}
Equivalently,
\[
\begin{aligned}
&\mathcal L(W_1^+,\dots,W_N^+) - \mathcal L(W_1,\dots,W_N)\\
&\qquad\leq
- \tfrac{\eta}{2} \sigma_{\min}(A)^{2(N-1)/N}
\|\nabla_A\mathcal F(A)\|_{\mathrm{HS}}^2.
\end{aligned}
\]
\end{theorem}

The proof, deferred to Appendix~\ref{sec:proof-linear-descent}, has two
parts. The geometric part establishes curvature-dependent smoothness of
the end-to-end pullback loss $\mathcal F$ on the tube
$\mathcal T_{R,D}$, with constant $L_{\kappa, R, D}$. The factorization
part expands the deep update $A^+ = W_N^+ \cdots W_1^+$ to first order in
$\eta$, applies the coercivity hypothesis to the leading-order term, and
absorbs the higher-order remainders into the step-size restriction. Under
iteration, this single-step inequality yields a linear rate
$\mathcal F(A(T)) - \mathcal F^\star \le e^{-c\, T}(\mathcal F(A(0)) -
\mathcal F^\star)$ provided the coercivity bound persists along the
trajectory; the latter is exactly the curvature-aware analogue of the
rank-margin persistence step of \cite{arora2018convergence}.

\paragraph{Where curvature enters.}
The balancedness term \(\sigma_{\min}(A)^{2(N-1)/N}\) is algebraic and
comes from the deep matrix factorization. The curvature of the output
manifold enters through the smoothness constant
\begin{equation}
\label{eq:LkRD}
L_{\kappa,R,D}
=
\frac1m\sum_{i=1}^m \|\xi_i\|^2
\left[
S_\kappa(R)^2 H_\kappa(D) + D\, B_\kappa(R)
\right],
\end{equation}
where \(S_\kappa(R) = \sinh(\sqrt{\kappa}R)/(\sqrt{\kappa}R)\) for
hyperbolic curvature \(-\kappa\). Negative curvature worsens the
allowable step size~\eqref{eq:linear-stepsize} through the expansion of
\(D\exp_{o_N}\), the Hessian of squared geodesic distance, and the second
derivative of the exponential map. The resulting per-step
contraction~\eqref{eq:linear-descent} is the same as in the Euclidean
case, but with a curvature-modulated step; Corollary~\ref{cor:curvature-stepsize}
below makes this curvature modulation quantitative.

\begin{corollary}[Curvature-explicit admissible step size]
\label{cor:curvature-stepsize}
Under the hypotheses of Theorem~\ref{thm:linear-descent}, the step-size
condition~\eqref{eq:linear-stepsize} is equivalent to
$0<\eta\le\eta^\star_\kappa$, where
\begin{equation}
\label{eq:eta-star-kappa}
\eta^\star_\kappa
:=
\frac{\sigma_{\min}(A)^{2(N-1)/N}}
     {4\,C_N\big(M^{3N-4}G_{\kappa,R,D}+L_{\kappa,R,D}M^{4N-4}\big)} .
\end{equation}
The constants $G_{\kappa,R,D}$ and $L_{\kappa,R,D}$ are non-decreasing in the
curvature magnitude $\kappa$: each grows through the angular-distortion
factor $S_\kappa(R)^2=\sinh^2(\sqrt{\kappa}R)/(\kappa R^2)$ and, in
$L_{\kappa,R,D}$ of~\eqref{eq:LkRD}, additionally through the
squared-distance Hessian bound $H_\kappa(D)$ and the exponential-map
second-derivative bound $B_\kappa(R)$. Consequently $\eta^\star_\kappa$ is
non-increasing in $\kappa$, equals its Euclidean value $\eta^\star_0$ at
$\kappa=0$, and in the curvature-dominated regime $\sqrt{\kappa}R\gg1$,
where $S_\kappa(R)^2\asymp e^{2\sqrt{\kappa}R}/(\kappa R^2)$ dominates both
constants, satisfies
\begin{equation}
\label{eq:eta-star-asymp}
\eta^\star_\kappa
=
O\!\Big(\tfrac{1}{S_\kappa(R)^2}\Big)
=
O\!\big(\kappa R^2 e^{-2\sqrt{\kappa}R}\big).
\end{equation}
Thus the admissible step size shrinks monotonically as curvature
increases---at least as fast as $1/S_\kappa(R)^2$ once the
$H_\kappa(D),B_\kappa(R)$ terms are included---and is exponentially small in
$\sqrt{\kappa}R$.
\end{corollary}

\begin{proof}
Condition~\eqref{eq:linear-stepsize} is affine in $\eta$ with strictly
positive coefficient $C_N(M^{3N-4}G_{\kappa,R,D}+L_{\kappa,R,D}M^{4N-4})$, so
for $\eta>0$ it is equivalent to $\eta\le\eta^\star_\kappa$. Write
$S_\kappa(R)=g(\sqrt{\kappa}R)$ with $g(t)=\sinh(t)/t=\sum_{k\ge0}t^{2k}/(2k+1)!$,
which is increasing on $[0,\infty)$; hence $S_\kappa(R)^2$ is increasing in
$\kappa$. The Hessian of the squared geodesic distance on the
hyperbolic factor has transverse eigenvalues $\sqrt{\kappa}\,r\coth(\sqrt{\kappa}\,r)$,
increasing in $\kappa$ for each $r\le D$, and the exponential-map
second-derivative bound is likewise non-decreasing in $\kappa$. Therefore
$G_{\kappa,R,D}$ and $L_{\kappa,R,D}$ are non-decreasing in $\kappa$, so
$\eta^\star_\kappa$ is non-increasing. Finally $\sinh^2(\sqrt{\kappa}R)\sim
\tfrac14 e^{2\sqrt{\kappa}R}$ as $\sqrt{\kappa}R\to\infty$
gives~\eqref{eq:eta-star-asymp}.
\end{proof}

Which of the two step-size regimes of Corollary~\ref{cor:curvature-stepsize}
is in force depends on the choice of loss. We make this precise in
Theorems~\ref{thm:tangent-surrogate} and~\ref{thm:intrinsic-rate} below: the
tangent-space surrogate is curvature-free, whereas the intrinsic loss is
curvature-modulated through $S_c(R)^2$ and collapses to the surrogate as
$c\to0^+$. Both statements require the Poincar\'e-ball machinery developed
next.


\paragraph{Hyperbolic specialization: Poincar\'e ball and M\"obius operations}
\label{sec:disc-hyperbolic}

For $c > 0$, the $n$-dimensional Poincar\'e ball is $\D_c^n := \{x \in
\R^n : c \|x\|^2 < 1\}$, equipped with Riemannian metric
$g_x^c = (\lambda_x^c)^2 g^E$ where $\lambda_x^c := 2/(1 - c\|x\|^2)$. The
induced hyperbolic distance is
$d_c(x, y) = \tfrac{2}{\sqrt c} \tanh^{-1}(\sqrt c\, \|{-x} \Mob y\|)$,
where $\Mob$ denotes M\"obius addition. The exponential and logarithmic
maps at the origin admit closed forms
$\exps(v) = \tanh(\sqrt c \|v\|)\, v/(\sqrt c \|v\|)$,
$\logs(y) = \tanh^{-1}(\sqrt c \|y\|)\, y/(\sqrt c \|y\|)$,
and are mutually inverse diffeomorphisms between $\D_c^n$ and $T_0\D_c^n
\cong \R^n$. For a Euclidean linear map $W : \R^n \to \R^m$, its
\emph{M\"obius version} is $W^{\Mobm}(x) := \exps(W \logs(x))$.
M\"obius matrix-vector multiplication is associative, yielding the
collapse identity
\begin{equation}
\label{eq:collapse}
\begin{aligned}
W_L \Mobm \cdots \Mobm W_1 \Mobm x &\;=\; \exps\!\big(W_{1:L}\, \logs(x)\big),\\
W_{1:L} &:= W_L W_{L-1} \cdots W_1.
\end{aligned}
\end{equation}
The deep linear hyperbolic neural network of
\cite{tifrea2018poincar} is precisely the instance of the general
layerwise Riemannian model of Section~\ref{sec:prelim-riem-nn} with
$\mathcal M_j = \D_c^{d_j}$ for all $j$.

\paragraph{Tangent-space surrogate loss: free reduction to Arora et al.}
\label{sec:disc-tangent}

Given data $\{(x^{(i)}, y^{(i)})\}_{i=1}^m \subset \D_c^{d_0} \times
\D_c^{d_L}$, write $\tilde x^{(i)} := \logs(x^{(i)})$, $\tilde y^{(i)} :=
\logs(y^{(i)})$ for the tangent-space lifts. The
\emph{tangent-space surrogate}
\(\mathcal E(W) := \tfrac{1}{2m} \sum_i \|W_{1:L}\tilde x^{(i)} -
\tilde y^{(i)}\|^2\) is, by the collapse identity~\eqref{eq:collapse} and
input whitening, exactly the Euclidean deep linear regression objective
$L^1(W_{1:L}) + \mathrm{const}$ on the tangent-lifted data $(\tilde
x^{(i)}, \tilde y^{(i)})$. The hyperbolic geometry contributes a change of
coordinates but no additional expressive power, and curvature enters only
through the dataset — not through the convergence mechanism.
Consequently, the discrete-time linear-rate theorem of
\cite{arora2018convergence} applies verbatim: under
$\delta$-balancedness with $\delta = (c')^2 / (256 L^3
\|\Phi\|_F^{2(L-1)/L})$, deficiency margin $c' > 0$, and step size
$\eta \leq (c')^{(4L-2)/L} / (6144 L^3 \|\Phi\|_F^{(6L-4)/L})$, the
iterates satisfy
\(\mathcal E(W(T)) \leq \varepsilon\) for every \(T \geq \tfrac{1}{\eta
(c')^{2(L-1)/L}} \log(\mathcal E(W(0))/\varepsilon)\).
We record this as a baseline. It follows from \cite{arora2018convergence}
by the change of variables above; the point we wish to isolate is that the
resulting rate carries no curvature factor.

\begin{theorem}[Tangent surrogate: curvature-free rate]
\label{thm:tangent-surrogate}
Consider the deep linear hyperbolic network $h(x)=\exps(W_{1:L}\logs(x))$ on
$\D_c^{d_0}\to\D_c^{d_L}$ trained on the tangent surrogate
$\mathcal E(W)=\tfrac1{2m}\sum_i\|W_{1:L}\tilde x^{(i)}-\tilde y^{(i)}\|^2$.
Assume the lifted regression is realizable: there is a target $\Phi$ with
$\Phi\tilde x^{(i)}=\tilde y^{(i)}$ for all $i$ (so $\min_W\mathcal E=0$), and
$\Phi$ has deficiency margin $c'>0$ (Definition~\ref{def:defmargin}). Run
gradient descent $W_j^{+}=W_j-\eta\nabla_{W_j}\mathcal E$ from a
$\delta$-balanced initialization (Definition~\ref{def:balanced}) with
\[
\begin{gathered}
\delta=\frac{(c')^2}{256\,L^3\|\Phi\|_F^{2(L-1)/L}},\\
0<\eta\le\frac{(c')^{(4L-2)/L}}{6144\,L^3\|\Phi\|_F^{(6L-4)/L}}.
\end{gathered}
\]
Write $\gamma:=\eta\,(c')^{2(L-1)/L}$ for the per-step rate. Then for every
$\varepsilon>0$ and every $T\ge\gamma^{-1}\log(\mathcal E(W(0))/\varepsilon)$,
the iterates satisfy $\mathcal E(W(T))\le\varepsilon$. The step size, the
contraction factor $1-\gamma$, and the iteration count are all independent of
the curvature $c$.
\end{theorem}

\begin{proof}
By the collapse identity~\eqref{eq:collapse}, $h(x^{(i)})=\exps(W_{1:L}\tilde
x^{(i)})$. With the lifts $\tilde x^{(i)}=\logs(x^{(i)})$, $\tilde
y^{(i)}=\logs(y^{(i)})$ and input whitening, the surrogate
\[
\mathcal E(W)=\tfrac1{2m}\sum_{i=1}^m\|W_{1:L}\tilde x^{(i)}-\tilde y^{(i)}\|^2
\]
is \emph{identically} the Euclidean deep-linear least-squares objective of
\cite{arora2018convergence} for the fixed data $\{(\tilde x^{(i)},\tilde
y^{(i)})\}_{i=1}^m$ with target $\Phi$. The maps $\exps,\logs$ enter only the
one-time construction of this data; they do not involve the optimization
variables $W_1,\dots,W_L$, and the loss landscape in $(W_1,\dots,W_L)$ is
exactly the Euclidean one. The discrete-time linear-convergence theorem of
\cite{arora2018convergence} therefore applies verbatim under the stated
$\delta$-balancedness, deficiency-margin, and step-size hypotheses, giving
$\mathcal E(W(T))\le\varepsilon$ for all $T$ as above, with per-step factor
$1-\gamma$. None of $\delta$, $c'$, $\eta$, or the contraction
factor depends on $c$: the curvature affects only the numerical values of
$\tilde x^{(i)},\tilde y^{(i)}$ (hence of $\Phi$ and of $\mathcal E(W(0))$),
not the convergence mechanism.
\end{proof}

\paragraph{Intrinsic squared hyperbolic distance loss}
\label{sec:disc-intrinsic}

The intrinsic loss \(\mathcal L_c(W) := \tfrac{1}{2m} \sum_i d_c
(\exps(W_{1:L}\tilde x^{(i)}), y^{(i)})^2\) is the geometrically correct
objective and does \emph{not} reduce to Euclidean linear regression after
$\logs$. Two obstructions distinguish this from the Euclidean case: the
pullback of a geodesically convex function along $\exps$ is generally not
Euclidean convex, and \(\mathcal L_c(W)\) is not convex in \(W\).
Convergence is nonetheless recovered through a \emph{bounded-ball metric
comparison}, which we now make precise.

\begin{lemma}[Bounded-ball metric comparison]
\label{lem:bounded-ball}
Let $c>0$, $R>0$, and $S_c(R):=\sinh(\sqrt c R)/(\sqrt c R)$. For all
$u,v\in T_{\mathbf 0}\D_c^n$ with $\|u\|,\|v\|\le R$,
\begin{equation}
\label{eq:bounded-ball}
\|u-v\|\;\le\;d_c(\exps u,\exps v)\;\le\;S_c(R)\,\|u-v\|,
\end{equation}
and consequently $\|u-v\|^2\le d_c(\exps u,\exps v)^2\le S_c(R)^2\|u-v\|^2$.
\end{lemma}

\begin{proof}
The Poincar\'e ball $\D_c^n$ has constant sectional curvature $-c$. We work in
geodesic normal coordinates at the origin $\mathbf 0$, so that
$T_{\mathbf 0}\D_c^n$ is identified with a Euclidean vector space and
$\exps:T_{\mathbf 0}\D_c^n\to \D_c^n$ denotes the Riemannian exponential map
at $\mathbf 0$. Throughout the proof, $\|\cdot\|$ denotes the norm on
$T_{\mathbf 0}\D_c^n$ induced by the Riemannian metric at $\mathbf 0$, while
$\|\cdot\|_g$ denotes the Riemannian norm on tangent spaces of $\D_c^n$.

Recall first the standard constant-curvature Jacobi-field computation. For
$w\in T_{\mathbf 0}\D_c^n$, the differential $d(\exps)_w$ acts as the identity
on the radial line $\mathbb R w$ and as multiplication by
\[
S_c(\|w\|):=\frac{\sinh(\sqrt c\,\|w\|)}{\sqrt c\,\|w\|}
\]
on the orthogonal complement $w^\perp$; see, for example,
\cite{boumal2023introduction}. Thus the singular values of $d(\exps)_w$ are
$1$ in the radial direction and $S_c(\|w\|)$ in the transverse directions.
In particular,
\[
\sigma_{\min}(d(\exps)_w)=1,
\qquad
\|d(\exps)_w\|_{\mathrm{op}}=S_c(\|w\|).
\]
Moreover $S_c(r)\ge 1$ and $S_c(r)$ is non-decreasing for $r\ge 0$.

We now prove the upper bound. The straight segment
\[
\sigma(t):=(1-t)u+tv,\qquad 0\le t\le 1,
\]
is a curve in the tangent space $T_{\mathbf 0}\D_c^n$, not a curve in the
Poincar\'e ball itself. The corresponding curve in the Poincar\'e ball is
\[
\gamma(t):=\exps(\sigma(t)).
\]
This curve joins the desired endpoints, since
\[
\gamma(0)=\exps u,
\qquad
\gamma(1)=\exps v.
\]
By definition, the Riemannian distance $d_c(p,q)$ is the infimum of the
Riemannian lengths of all piecewise smooth curves joining $p$ to $q$. Hence,
for this particular admissible curve $\gamma$,
\[
d_c(\exps u,\exps v)\le L_g(\gamma),
\]
where
\[
L_g(\gamma)=\int_0^1 \|\dot\gamma(t)\|_g\,dt.
\]
Using the chain rule,
\[
\dot\gamma(t)=d(\exps)_{\sigma(t)}\,\dot\sigma(t).
\]
Therefore
\[
d_c(\exps u,\exps v)
\le
\int_0^1
\big\|d(\exps)_{\sigma(t)}\,\dot\sigma(t)\big\|_g\,dt.
\]
Since $\sigma(t)$ is the Euclidean segment between $u$ and $v$, convexity of
the norm gives
\[
\|\sigma(t)\|
\le
(1-t)\|u\|+t\|v\|
\le
R.
\]
Thus, for every $t\in[0,1]$,
\[
\|d(\exps)_{\sigma(t)}\|_{\mathrm{op}}
=
S_c(\|\sigma(t)\|)
\le
S_c(R).
\]
It follows that
\[
\begin{aligned}
d_c(\exps u,\exps v)
&\le
\int_0^1
\big\|d(\exps)_{\sigma(t)}\,\dot\sigma(t)\big\|_g\,dt \\
&\le
S_c(R)\int_0^1 \|\dot\sigma(t)\|\,dt.
\end{aligned}
\]
Finally, $\dot\sigma(t)=v-u$ is constant, so
\[
\int_0^1\|\dot\sigma(t)\|\,dt
=
\|v-u\|.
\]
Therefore
\[
d_c(\exps u,\exps v)
\le
S_c(R)\|u-v\|.
\]

We next prove the lower bound. Since $\exps$ is a diffeomorphism from
$T_{\mathbf 0}\D_c^n$ onto $\D_c^n$, its inverse is the logarithmic map
$\logs:\D_c^n\to T_{\mathbf 0}\D_c^n$. For $a\in\D_c^n$,
\[
d(\logs)_a
=
\bigl(d(\exps)_{\logs a}\bigr)^{-1}.
\]
The singular values of $d(\exps)_{\logs a}$ are $1$ and
$S_c(\|\logs a\|)$, so the singular values of $d(\logs)_a$ are
$1$ and
\[
\frac{1}{S_c(\|\logs a\|)}\le 1.
\]
Hence
\[
\|d(\logs)_a\|_{\mathrm{op}}\le 1
\]
for every $a\in\D_c^n$. Equivalently, infinitesimal lengths do not increase
under $\logs$.

To turn this infinitesimal estimate into a distance estimate, let
$\eta:[0,1]\to\D_c^n$ be any piecewise smooth curve joining
$\exps u$ to $\exps v$. Then $\logs\circ\eta$ is a curve in
$T_{\mathbf 0}\D_c^n$ joining $u$ to $v$, and
\[
\frac{d}{dt}(\logs\circ\eta)(t)
=
d(\logs)_{\eta(t)}\,\dot\eta(t).
\]
Therefore
\[
\begin{aligned}
L(\logs\circ\eta)
&=
\int_0^1
\big\|d(\logs)_{\eta(t)}\,\dot\eta(t)\big\|\,dt \\
&\le
\int_0^1
\|\dot\eta(t)\|_g\,dt
=
L_g(\eta).
\end{aligned}
\]
Since the Euclidean distance between the endpoints of a curve is at most the
length of the curve,
\[
\|u-v\|
\le
L(\logs\circ\eta)
\le
L_g(\eta).
\]
This holds for every piecewise smooth curve $\eta$ joining $\exps u$ to
$\exps v$. Taking the infimum over all such curves gives
\[
\|u-v\|
\le
d_c(\exps u,\exps v).
\]

Combining the two estimates yields
\[
\|u-v\|
\le
d_c(\exps u,\exps v)
\le
S_c(R)\|u-v\|.
\]
Squaring both sides gives
\[
\|u-v\|^2
\le
d_c(\exps u,\exps v)^2
\le
S_c(R)^2\|u-v\|^2,
\]
as claimed.
\end{proof}
The lower bound in~\eqref{eq:bounded-ball} expresses that $\exps$ is distance
non-decreasing in nonpositive curvature, and the upper bound is the Rauch
expansion factor $S_c(R)$. The intrinsic loss is thereby sandwiched between
the surrogate and an $S_c(R)^2$-inflated surrogate, and this is the mechanism
that transfers the tangent-loss convergence of
Theorem~\ref{thm:tangent-surrogate} to $\mathcal L_c$ at a curvature-dependent
price.

\begin{theorem}[Intrinsic loss: curvature-modulated rate]
\label{thm:intrinsic-rate}
Adopt the hypotheses of Theorem~\ref{thm:tangent-surrogate} (in particular
realizability, so $\min_W\mathcal E=\min_W\mathcal L_c=0$), and let
$\gamma=\eta(c')^{2(L-1)/L}$ be as there. Fix $R>0$ and suppose the
gradient-descent iterates of Theorem~\ref{thm:tangent-surrogate} obey the
bounded-radius condition $\|W_{1:L}(t)\tilde x^{(i)}\|\le R$ and
$\|\tilde y^{(i)}\|\le R$ for all $t\ge0$ and all $i$. Then the intrinsic
loss along the same iterates satisfies, for every $T\ge0$,
\begin{equation}
\label{eq:intrinsic-rate}
\mathcal L_c(W(T))\;\le\;S_c(R)^2\,e^{-\gamma T}\,\mathcal L_c(W(0)).
\end{equation}
In particular $\mathcal L_c(W(T))\to0$ at a linear rate; the curvature penalty
is the prefactor $S_c(R)^2=\big(\sinh(\sqrt c R)/(\sqrt c R)\big)^2\ge1$; and
as $c\to0^+$, $S_c(R)^2\to1$, so~\eqref{eq:intrinsic-rate} collapses to the
curvature-free rate of Theorem~\ref{thm:tangent-surrogate}.
\end{theorem}

\begin{proof}
Write $r_i(t):=W_{1:L}(t)\tilde x^{(i)}-\tilde y^{(i)}$ for the $i$-th tangent
residual. Apply Lemma~\ref{lem:bounded-ball} with $u=W_{1:L}(t)\tilde x^{(i)}$
and $v=\tilde y^{(i)}$, both in the ball of radius $R$ by hypothesis; since
$\exps(\tilde y^{(i)})=y^{(i)}$ and
$\exps(W_{1:L}(t)\tilde x^{(i)})=h_t(x^{(i)})$,
\[
\begin{aligned}
\|r_i(t)\|^2
&\le d_c\big(h_t(x^{(i)}),y^{(i)}\big)^2\\
&\le S_c(R)^2\,\|r_i(t)\|^2 .
\end{aligned}
\]
Averaging over $i$ and dividing by $2$ gives the loss sandwich
\begin{equation}
\label{eq:loss-sandwich}
\mathcal E(W)\;\le\;\mathcal L_c(W)\;\le\;S_c(R)^2\,\mathcal E(W),
\end{equation}
valid for every $W$ on the trajectory. Realizability gives
$\min_W\mathcal E=0$ at $W_{1:L}=\Phi$, where $r_i\equiv0$ forces
$d_c(h(x^{(i)}),y^{(i)})=0$, so $\min_W\mathcal L_c=0$ as well. By
Theorem~\ref{thm:tangent-surrogate},
$\mathcal E(W(T))\le e^{-\gamma T}\mathcal E(W(0))$, so using the upper side
of~\eqref{eq:loss-sandwich} at time $T$ and the lower side at time $0$,
\[
\begin{aligned}
\mathcal L_c(W(T))
&\le S_c(R)^2\,\mathcal E(W(T))\\
&\le S_c(R)^2 e^{-\gamma T}\,\mathcal E(W(0))\\
&\le S_c(R)^2 e^{-\gamma T}\,\mathcal L_c(W(0)),
\end{aligned}
\]
which is~\eqref{eq:intrinsic-rate}. Finally
$S_c(R)=\sinh(\sqrt c R)/(\sqrt c R)=1+cR^2/6+O(c^2)\to1$ as $c\to0^+$, so the
prefactor tends to $1$ and, by~\eqref{eq:loss-sandwich},
$\mathcal L_c\to\mathcal E$ uniformly on the bounded region.
\end{proof}

\begin{remark}[Direct minimization of $\mathcal L_c$]
\label{rem:direct-intrinsic}
Theorem~\ref{thm:intrinsic-rate} runs gradient descent on the surrogate
$\mathcal E$ and transfers the guarantee to $\mathcal L_c$, so the curvature
cost surfaces as the prefactor $S_c(R)^2$. Minimizing $\mathcal L_c$ directly
is expected instead to trade this prefactor for a curvature-slowed exponent,
with per-step contraction $1-\gamma/S_c(R)^2$: the lower bound
in~\eqref{eq:bounded-ball}, together with the value
sandwich~\eqref{eq:loss-sandwich}, yields a Polyak--\L ojasiewicz inequality
for $\mathcal L_c$ with constant degraded by $S_c(R)^2$ on bounded sublevel
sets. Establishing the matching smoothness constant for $\mathcal L_c$ along
the trajectory, and the continuous-time gradient-flow analogue, is left to
future work.
\end{remark}


% =====================================================================
\subsection{Nonlinear case}
\label{sec:main-nonlinear}
% =====================================================================

We now consider the general case where the tangent-space operations
$F_j$ may be nonlinear. The same recursive layerwise constants of
Section~\ref{sec:prelim-recursive} control optimization (the descent
lemma, Theorem~\ref{thm:general-descent}) and generalization (the
covering-number bound, Theorem~\ref{thm:layerwise-generalization}). The
linear case of Section~\ref{sec:main-linear} is recovered by
specializing $F_j$ to a Euclidean linear map.

Recall from Section~\ref{sec:prelim-recursive} the recursive smoothness
constants $J_j, K_j$ of~\eqref{eq:JK-recursive} and the network smoothness
constant $L_{\mathrm{net}}$ of~\eqref{eq:Lnet-def} defined for a general
layerwise Riemannian neural network. Our first main result is a
single-step descent inequality whose admissible step size is controlled
by $L_{\mathrm{net}}$ and which therefore inherits a curvature-dependent
restriction through the constants $\alpha_j,\beta_j,\gamma_j,\delta_j$ of
the layerwise exponential and logarithmic maps.

\begin{theorem}[General layerwise Riemannian descent lemma]
\label{thm:general-descent}
Assume that the parameter trajectory remains in the layerwise geometric
tube of Section~\ref{sec:prelim-riem-nn}, so that \(h_j^{(i)}(\theta)\in
\mathcal T_j\) for all \(i,j\). Assume also that
\(\mathcal L:\Theta_1\times\cdots\times\Theta_N\to\mathbb R\) is
differentiable and that the constants $\alpha_j, \beta_j, \gamma_j,
\delta_j$, $P_j, Q_j, R_j^\theta, S_j^\theta, T_j^\theta$, $G_N, H_N$ above
are finite. Let \(\theta^+ = \theta-\eta\nabla_\theta\mathcal L(\theta)\).
If the segment \(\theta_s:=\theta-s\eta\nabla_\theta\mathcal L(\theta)\),
\(s\in[0,1]\), stays inside the parameter tube, and if
\begin{equation}
\label{eq:general-stepsize}
0<\eta\le \frac1{L_{\mathrm{net}}},
\end{equation}
then
\begin{equation}
\label{eq:general-descent}
\mathcal L(\theta^+)
\le
\mathcal L(\theta)
-
\frac{\eta}{2}
\|\nabla_\theta\mathcal L(\theta)\|^2.
\end{equation}
\end{theorem}

The proof is a curvature-aware smoothness argument that propagates the
recursive bounds of Section~\ref{sec:prelim-recursive} through the chain
rule for $\mathcal L$ to obtain \(L_{\mathrm{net}}\)-smoothness on the
tube; the full argument is given in
Appendix~\ref{sec:proof-general-descent}.

\paragraph{Where each manifold curvature enters.}
The curvature of \(\mathcal M_j\) enters through
\(\alpha_j = \sup \|D\exp_{o_j}^{(j)}\|\),
\(\gamma_j = \sup \|D^2\exp_{o_j}^{(j)}\|\),
\(\beta_j = \sup \|D\log_{o_j}^{(j)}\|\), and
\(\delta_j = \sup \|D^2\log_{o_j}^{(j)}\|\). Thus the step size
restriction~\eqref{eq:general-stepsize}, with $L_{\mathrm{net}}$ given
by~\eqref{eq:Lnet-def}, depends on all layerwise
curvatures \(\kappa_0,\kappa_1,\dots,\kappa_N\). For hyperbolic manifolds,
\(\alpha_j \lesssim S_{\kappa_j}(R_j) = \sinh(\sqrt{\kappa_j}R_j)/(\sqrt{\kappa_j}R_j)\),
so more negative curvature or larger radius \(R_j\) worsens the smoothness
constant and reduces the admissible step size.

We now turn to generalization. The closest manifold baseline is
\cite{sarkar2025learning}, who give a curvature-adaptive Rademacher bound
for a single $L$-Lipschitz function on a Riemannian manifold $\mathcal M_0$
via a Bishop--Gromov covering of $\mathcal M_0$. In that bound the
network is treated as a black-box $L$-Lipschitz function on $\mathcal M_0$,
so the curvature of any \emph{hidden} representation is invisible. Our
generalization bound keeps the curvature-adaptive covering of
\cite{sarkar2025learning}, but replaces the monolithic Lipschitz constant
$L$ by a \emph{layerwise} expression built from the same recursive
constants $\alpha_j, \beta_j, P_j$ that drove
Theorem~\ref{thm:general-descent} above; thus optimization and
generalization for the nonlinear model live in one set of layerwise
geometric constants.

\paragraph{Network spatial Lipschitz constant.}
For the model of Section~\ref{sec:prelim-riem-nn}, define recursively
\begin{equation}
\label{eq:Lambda-recursive}
\Lambda_0 := 1,
\qquad
\Lambda_j := \alpha_j P_j \beta_{j-1} \Lambda_{j-1}
\quad (j=1,\dots,N),
\end{equation}
so that $\Lambda_N = \prod_{j=1}^N \alpha_j P_j \beta_{j-1}$. As shown in
the proof, $h_N(\cdot;\theta) : \mathcal M_0 \to \mathcal M_N$ is
$\Lambda_N$-Lipschitz with respect to the geodesic distances
$d_{\mathcal M_0}, d_{\mathcal M_N}$ for every $\theta$ in the parameter
tube. The curvature of every hidden manifold $\mathcal M_j$ enters
$\Lambda_N$ multiplicatively through $\alpha_j$; e.g., for hyperbolic
$\mathcal M_j = \mathbb H^{d_j}_{-\kappa_j}$,
$\alpha_j \lesssim S_{\kappa_j}(R_j)$, so layerwise curvature
contributions are additive in $\log \Lambda_N$.

\begin{theorem}[Curvature-adaptive generalization bound for layerwise Riemannian neural networks]
\label{thm:layerwise-generalization}
Let $\mathcal M_0$ be a compact $d_0$-dimensional Riemannian manifold
with sectional curvature bounded by
$\kappa_{\min}^{(0)} \le \mathrm{sec} \le \kappa_{\max}^{(0)}$, injectivity
radius $\iota_0 > 0$, and finite volume. Assume the general layerwise
Riemannian neural network of Section~\ref{sec:prelim-riem-nn} with
parameters in a compact tube and finite layerwise constants
$\alpha_j, \beta_j, P_j$. Assume the loss
$\ell : \mathcal M_N \to \mathbb R$ is $L_\ell$-Lipschitz on $\mathcal T_N$
and bounded by $B$. Let $\{x_i\}_{i=1}^n$ be drawn i.i.d.\ from a
probability measure on $\mathcal M_0$, and let
$\hat R_n(\theta) := \tfrac1n \sum_{i=1}^n \ell(h_N(x_i;\theta))$ and
$R(\theta) := \mathbb E[\ell(h_N(x;\theta))]$. Then for any
$\delta \in (0,1)$, with probability at least $1-\delta$, for all
$\theta$ in the parameter tube,
\begin{equation}
\label{eq:layerwise-gen}
\begin{aligned}
|R(\theta) - \hat R_n(\theta)|
\le
O\Bigg(
& L_\ell \cdot
\frac{\sqrt{d_0 \log(L_\ell \Lambda_N \sqrt n)}}{n^{1/d_0}}\\
& + L_\ell \cdot
\frac{\psi(\kappa_{\max}^{(0)}, L_\ell \Lambda_N)}{n^{1/d_0}}\\
& + B \sqrt{\frac{\log(1/\delta)}{n}}
\Bigg),
\end{aligned}
\end{equation}
where $\Lambda_N = \prod_{j=1}^N \alpha_j P_j \beta_{j-1}$ is the network
spatial Lipschitz constant and $\psi(\kappa, L) := \sqrt{|\kappa|}/L$ is
the Bishop--Gromov curvature penalty of \cite{sarkar2025learning}.
\end{theorem}

The proof, given in Appendix~\ref{sec:proof-layerwise-gen}, is a
layerwise refinement of \cite{sarkar2025learning}: the layerwise
Jacobian and log/exp bounds of Section~\ref{sec:prelim-recursive}
collapse the network to a single $L_\ell \Lambda_N$-Lipschitz function on
$\mathcal M_0$, after which the Bishop--Gromov covering argument and
Dudley entropy integral of \cite{sarkar2025learning} apply. Standard
symmetrization and bounded-difference concentration give the explicit
high-probability statement.

\paragraph{Comparison with \cite{sarkar2025learning}.}
\cite{sarkar2025learning} treats the network as a single $L$-Lipschitz
function on $\mathcal M_0$; the curvature of any hidden representation
is invisible. Theorem~\ref{thm:layerwise-generalization} replaces the
monolithic $L$ by the layerwise product
$\Lambda_N$ of~\eqref{eq:Lambda-recursive}. This is sharper whenever
the layerwise constants are individually controlled: for a depth-$N$
hyperbolic network with bounded radius $R_j$ at each layer,
$\log \Lambda_N \le \sum_j \log(\alpha_j P_j \beta_{j-1})$, and
$\alpha_j \lesssim S_{\kappa_j}(R_j)$ contributes
$\log S_{\kappa_j}(R_j) \approx \kappa_j R_j^2 / 6$ in the small-curvature
regime and $\sqrt{\kappa_j} R_j$ for large negative curvature. When
$\mathcal M_0$ is Euclidean and $\beta_j = \alpha_j = 1$ for every $j$,
\(\psi \to 0\) and~\eqref{eq:layerwise-gen} reduces to the standard
$L_\ell \Lambda_N$ Rademacher bound for Lipschitz function classes.

\paragraph{Where each manifold curvature enters.}
In~\eqref{eq:layerwise-gen}, the curvature of the \emph{input} manifold
$\mathcal M_0$ enters through the Bishop--Gromov term
$\psi(\kappa_{\max}^{(0)}, L_\ell \Lambda_N)$: it governs the volume
growth of geodesic balls and hence the covering number of $\mathcal M_0$.
The curvature of every \emph{hidden} manifold $\mathcal M_j$ ($j \ge 1$)
enters through $\Lambda_N$ in~\eqref{eq:Lambda-recursive} and hence in
the $\log(L_\ell \Lambda_N \sqrt n)$ term of~\eqref{eq:layerwise-gen}:
it inflates the Jacobians $\alpha_j$ of $\exp_{o_j}^{(j)}$, and hence the
network's overall spatial Lipschitz constant, but does not change the
covering geometry of $\mathcal M_0$ itself. The two contributions are
quantitatively different: input-manifold curvature gives a
$\sqrt{|\kappa_{\max}^{(0)}|}$ contribution, hidden-layer curvature a
$\log S_\kappa(R)$ contribution.


% =====================================================================
\subsection{Graph as a special case}
\label{sec:main-graph}
% =====================================================================

As a downstream application of the nonlinear case
(Theorem~\ref{thm:layerwise-generalization}), we specialize to the
pseudo-Riemannian Q-GCN of \cite{xiong2022pseudo}, recovering a
curvature-adaptive transductive generalization bound for graph networks.

The pseudo-Riemannian Q-GCN of Section~\ref{sec:prelim-qgcn} is not a
direct instance of Theorem~\ref{thm:layerwise-generalization}, because
the pseudo-hyperboloid $\mathcal Q_\beta^{s,t}$ is not geodesically
complete and the standard learning setting is replaced by a transductive
one. Both gaps are bridged by minor adjustments: \textup{(i)} cover the
diffeomorphic product manifold
$M_\psi = \prod_\ell (\mathbb S^t_{-\beta_\ell} \times \mathbb R^{s_\ell})$
through the spherical projection $\psi$ of \cite{xiong2022pseudo}, so
that the Bishop--Gromov step is applied factorwise; \textup{(ii)} replace
the standard PAC concentration step of
Theorem~\ref{thm:layerwise-generalization} by the transductive Rademacher
comparison of \cite[Theorem~4.3]{tang2023towards}. Step \textup{(i)} is
exactly the reason an additive layerwise curvature penalty
$\Psi_{\mathrm{curv}} = \sum_\ell c_\ell \sqrt{|\kappa^{(\ell)}_{\max}|}$
appears in the Q-GCN case: each spherical factor of $M_\psi$ contributes
its own Bishop--Gromov term.

\begin{corollary}[Curvature-adaptive transductive generalization bound for Q-GCN]
\label{thm:prgcn_transductive_generalization}
Let $\mathcal Q_\beta^{s,t}$ be the pseudo-hyperboloid used by Q-GCN with
spherical projection $\psi$, and let the reference point $o$ be chosen so
that the induced tangent-space operations are identified with the tangent
space of the pseudo-hyperboloid. Assume the transductive learning setting
and stochastic-optimization assumptions of
\cite[Theorem~4.3]{tang2023towards} hold, the Q-GCN layer is as in
Section~\ref{sec:prelim-qgcn}, all maps used in that layer are
well-defined on the forward trajectory, and the node-wise loss class is
uniformly bounded by $B$ and $L_Q$-Lipschitz with respect to the
product-geodesic metric on $\mathbb S^t_{-\beta} \times \R^s$. Let
$\kappa_{\max}^{(\ell)}$ be an upper sectional-curvature bound for the
spherical factor at layer $\ell$ and define
$\Psi_{\mathrm{curv}} := \sum_\ell c_\ell \sqrt{|\kappa_{\max}^{(\ell)}|}$.
Then for any $\delta \in (0,1)$ and the learning rate schedule of
\cite[Theorem~4.3]{tang2023towards}, with probability $\geq 1 - \delta$,
\begin{equation}
\label{eq:qgcn-gen}
\begin{aligned}
& R_u(w^{(T+1)}) - R_m(w^{(T+1)})\\
& \quad\;=\; \widetilde O\!\left(
\tfrac{(m+u)^{3/2}}{mu} \cdot T^*_\alpha(T) \cdot \log\tfrac1\delta \cdot
\mathfrak C_Q(T, m+u) \right),
\end{aligned}
\end{equation}
where $T^*_\alpha(T)$ is the optimization-radius factor of
\cite{tang2023towards} (with the three regimes of $\alpha$) and
\begin{equation}
\label{eq:CQ-def}
\begin{aligned}
\mathfrak C_Q(T, n) \;:=\; L_Q
\Big[ & d_{\mathrm{eff}} \log\!\big(L_Q n^{1/d_{\mathrm{eff}}}\big)
+ \Psi_{\mathrm{curv}}\\
& + \log\!\big(\mathrm{Vol}(M_\psi)/\omega_{d_{\mathrm{eff}}}\big)\Big]^{1/2}.
\end{aligned}
\end{equation}
\end{corollary}

The proof, given in Appendix~\ref{sec:proof-gen}, applies
Theorem~\ref{thm:layerwise-generalization} after the diffeomorphic chart
and transductive replacements above. The Euclidean architecture constant
$L_F$ of \cite{tang2023towards} is replaced by
$\mathfrak C_Q(T, m+u)$ of~\eqref{eq:CQ-def}, which factorizes into three
contributions: an effective-dimension entropy
$d_{\mathrm{eff}} \log(\cdots)$, the layerwise curvature penalty
$\Psi_{\mathrm{curv}}$, and a volume term
$\log(\mathrm{Vol}(M_\psi)/\omega_{d_{\mathrm{eff}}})$. As all
$\kappa_{\max}^{(\ell)} \to 0$ the curvature penalty vanishes
and~\eqref{eq:qgcn-gen} reduces to the Euclidean expression of
\cite{tang2023towards}.


\paragraph{Pseudo-Riemannian setting: spacelike trajectory hypothesis}
\label{sec:disc-pseudo}

To lift the hyperbolic specialization to the full pseudo-Riemannian
setting on $Q_\beta^{s,t}$, two issues arise. First, the diffeomorphic
$\widehat{\log}^\beta_o, \widehat{\exp}^\beta_o$ replace
$\logs, \exps$ in the model, but the same collapse identity holds under
uniform curvature: $x_L = \widehat{\exp}^\beta_o(W_{1:L}
\widehat{\log}^\beta_o(x_0))$. The tangent-space surrogate of
Section~\ref{sec:disc-tangent} therefore lifts verbatim. Second, the
intrinsic pseudo-Riemannian squared distance $d_\beta^2$ takes positive
values along spacelike geodesics, negative values along timelike
geodesics, and zero along null geodesics; it is only well-defined as a
non-negative training objective in the \emph{spacelike sector}. Under a
working hypothesis that the gradient-descent trajectory stays in the
spacelike sector, the spacelike-restricted metric has the same polar-form
$d\tau^2 + (\sqrt{-\beta}\sinh(\tau/\sqrt{-\beta}))^2 d\omega^2$ as the
hyperbolic metric of curvature $-c$ with $c = 1/|\beta|$, and the
hyperbolic argument of Section~\ref{sec:disc-intrinsic} carries through
with $S_c(R)$ replaced by $S_\beta(R) := \sinh(R/\sqrt{-\beta}) /
(R/\sqrt{-\beta})$. An alternative formulation — using the
$\psi$-Riemannian product metric pulled back from $\mathbb S^t_{-\beta}
\times \R^s$ — avoids the spacelike-trajectory hypothesis at the cost of
departing from the genuine indefinite-metric formulation.


% =====================================================================
\section{Future Directions}
\label{sec:future}
% =====================================================================


Three further extensions are natural:
\textup{(i)} \emph{M\"obius biases.} The bias-augmented model $x \mapsto W
\otimes_c x \oplus_c b$ breaks the collapse identity~\eqref{eq:collapse}
and prevents end-to-end parameterization by a single matrix; a
perturbative extension treating biases as $O(\beta)$ perturbations of the
bias-free flow is open. Riemannian gradient descent on the biases is a
natural alternative.
\textup{(ii)} \emph{Nonlinear activations.} Our analysis, like
\cite{arora2018convergence}, requires linearity. Extending to hyperbolic
ReLU or hyperbolic sigmoid (pointwise nonlinearities via
$\varphi^{\Mobm}$) is open even in the Euclidean case and provably harder
in ours because the non-convexity of the intrinsic loss compounds with
the non-convexity introduced by the activation.
\textup{(iii)} \emph{Tightening curvature dependence of the generalization
bound.} Our $\Psi_{\mathrm{curv}}$ is layerwise additive in
$\sqrt{|\kappa_{\max}^{(\ell)}|}$, but lower bounds in this geometric
regime are unknown.

% =====================================================================
\section{Experiments}
\label{sec:experiments}
% =====================================================================

We validate the optimization side of our theory --- the descent lemma
(Theorem~\ref{thm:general-descent}), its deep-linear specialization
(Theorem~\ref{thm:linear-descent}), and the intrinsic-loss specialization
of Section~\ref{sec:disc-intrinsic} --- on the deep-linear hyperbolic
network of Section~\ref{sec:disc-hyperbolic}, i.e.\ the instance of the
general layerwise model (Section~\ref{sec:prelim-riem-nn}) with
$M_j = \D^{d_j}_c$ for every layer. All experiments probe the way a single
scalar, the angular distortion factor
$S_c(R) = \sinh(\sqrt c R)/(\sqrt c R)$, governs optimization; this is the
factor through which curvature enters both $L_{\mathrm{net}}$ and the
bounded-ball comparison of Section~\ref{sec:disc-intrinsic}.

\paragraph{Setup.}
We implement the closed-form Poincar\'e-ball maps $\exps,\logs$ of
Section~\ref{sec:disc-hyperbolic}, the collapse
identity~\eqref{eq:collapse} $h = \exps(W_{1:L}\logs(x))$, the
tangent-space surrogate $\mathcal E$ of Section~\ref{sec:disc-tangent},
and the intrinsic squared-distance loss $\mathcal L_c$ of
Section~\ref{sec:disc-intrinsic}. Factor gradients are computed
analytically: we form the end-to-end gradient $G = \nabla_A \mathcal L_c$
in closed form using~\eqref{eq:pullback-gradient}, then distribute it to
the factors by the identity
$\nabla_{W_j}\mathcal L_c = W_{j+1:N}^{*}\, G\, W_{1:j-1}^{*}$ established
in Step~3 of the proof of Theorem~\ref{thm:linear-descent}. The analytic
gradient agrees with central finite differences to relative error
$\sim\!10^{-9}$ across $c \in \{0,\tfrac12,1,4\}$, which also serves as a
numerical check of that identity. Unless stated otherwise we use depth
$N=3$, width $d=6$, $m=300$ samples, bounded radius $R=0.8$, and an
initialization that is exactly $\delta$-balanced at $t=0$ (each $W_j$ a
scaled orthogonal matrix, so the defect of
Definition~\ref{def:balanced} vanishes initially). A single
tangent-space regression problem is generated once and reused at every
curvature via $x^{(i)} = \exps(\tilde x^{(i)})$,
$y^{(i)} = \exps(\tilde y^{(i)})$, so that curvature enters
\emph{only} through the geometry and never through the data; this is the
experimental counterpart of holding the tangent-lifted problem fixed in
the comparison $\mathcal E \le \mathcal L_c \le S_c(R)^2\mathcal E$.

\paragraph{What the theory predicts, and what it does not.}
The bounded-ball comparison of Section~\ref{sec:disc-intrinsic} couples
$\mathcal L_c$ to the Euclidean surrogate $\mathcal E$ through
$S_c(R)^2$. Because that factor appears in the \emph{gradient} comparison
$\|\nabla\mathcal L_c\|_F \le S_c(R)^2 \|\nabla\mathcal E\|_F$, it acts as
a curvature-dependent amplifier of the effective gradient. Its theoretical
consequence is therefore a restriction on the admissible step size: the
larger $S_c(R)$, the smaller the $\eta$ for which the descent guarantee
of Theorem~\ref{thm:general-descent} (step-size
condition~\eqref{eq:general-stepsize}) and the curvature-dependent
condition~\eqref{eq:linear-stepsize} of
Theorem~\ref{thm:linear-descent} can hold. We emphasize that the theory does \emph{not} predict that the raw
per-step contraction \emph{rate} degrades monotonically with curvature at
a fixed step size; under a fixed $\eta$ the same amplification accelerates
the iteration until it destabilizes. The clean monotone signature of
curvature is thus expected in the maximum stable step size, not in the
fixed-step rate, and our experiments are organized accordingly.

\begin{figure*}[t]
\centering
\includegraphics[width=\textwidth]{figure_descent.pdf}
\caption{Optimization of the deep-linear hyperbolic network of
Section~\ref{sec:disc-hyperbolic} as a function of curvature $c$ at fixed
radius $R=0.8$. \textbf{(a)} Loss-excess trajectories at fixed step size on
a single fixed tangent-space problem: the intrinsic-loss curve at $c=0$
coincides with the Euclidean tangent-surrogate curve (dashed), the $c\to0$
collapse of Section~\ref{sec:disc-intrinsic}. \textbf{(b)} Maximum stable
step size $\eta^\star$ decreases monotonically with curvature, spanning
roughly two orders of magnitude; this is the empirical signature
of~\eqref{eq:general-stepsize} with $L_{\mathrm{net}}$
of~\eqref{eq:Lnet-def} increasing in curvature
(Theorems~\ref{thm:general-descent} and~\ref{thm:linear-descent}).
\textbf{(c)} The same $\eta^\star$, normalized to its $c=0$ value, falls
\emph{faster} than the curvature factor $1/S_c(R)^2$ alone, consistent
with $L_{\kappa,R,D}$ of~\eqref{eq:LkRD} carrying the additional Hessian
and exponential-map second-derivative terms $H_\kappa(D), B_\kappa(R)$
(Theorem~\ref{thm:linear-descent}). \textbf{(d)} The
$\delta$-balancedness defect
$\max_j\|W_{j+1}^\top W_{j+1} - W_j W_j^\top\|_F$ remains small
($\lesssim 3\times10^{-3}$) along the entire trajectory at every
curvature, so the structural hypothesis of
Theorem~\ref{thm:linear-descent} is maintained during training.}
\label{fig:descent}
\end{figure*}

\paragraph{(a) Convergence and the $c\to0$ collapse.}
Figure~\ref{fig:descent}(a) shows loss-excess trajectories under
full-batch gradient descent at a fixed step size, on the single fixed
tangent-space problem, for $c \in \{0,\tfrac12,1,2,4,8\}$, together with
the Euclidean tangent-surrogate trajectory. At $c=0$ the intrinsic-loss
trajectory is indistinguishable from the tangent-surrogate trajectory:
this is the collapse $\mathcal L_c \to \mathcal E$ as $c \to 0^+$ asserted
in Section~\ref{sec:disc-intrinsic}, observed to machine precision. For
$c>0$ the trajectories separate, but --- as anticipated above --- the
fixed-step iteration is \emph{faster}, not slower, at higher curvature,
because $S_c(R)^2$ amplifies the effective gradient. This panel is a
qualitative sanity check of the collapse and of the gradient-amplification
mechanism, not a rate law.

\paragraph{(b) Admissible step size shrinks with curvature.}
Figure~\ref{fig:descent}(b) reports the maximum stable step size
$\eta^\star$, defined as the largest $\eta$ on a geometric grid for which
the iteration remains bounded and decreases the loss over a fixed horizon.
$\eta^\star$ is strictly monotonically decreasing in $c$, dropping by
roughly a factor of $80$ from $c=0$ to $c=8$. This is the central
empirical result: it is the direct signature of the curvature-dependent
step-size restriction~\eqref{eq:general-stepsize} of
Theorem~\ref{thm:general-descent}, with $L_{\mathrm{net}}$
of~\eqref{eq:Lnet-def} increasing in curvature through
$\alpha_j,\gamma_j$, and of condition~\eqref{eq:linear-stepsize} of
Theorem~\ref{thm:linear-descent}. Because $L_{\mathrm{net}}$ is an upper
bound, we report the monotone ordering rather than a match of constants.

\paragraph{(c) Beyond the bare $S_c(R)^2$ factor.}
Figure~\ref{fig:descent}(c) normalizes $\eta^\star$ to its $c=0$ value and
compares it against $1/S_c(R)^2$. The measured step size falls faster than
$1/S_c(R)^2$ alone. This is consistent with, and indeed evidence for, the
full curvature-dependent smoothness constant $L_{\kappa,R,D}$
of~\eqref{eq:LkRD}: there curvature enters not only through
$S_\kappa(R)^2$ but also through the squared-distance Hessian bound
$H_\kappa(D)$ and the exponential-map second-derivative bound
$B_\kappa(R)$. A model in which only the angular-distortion factor
contributed would \emph{under}-predict the step-size penalty; the data sit
below that prediction, indicating that the additional terms in
$L_{\kappa,R,D}$ are doing real work.

\paragraph{(d) Balancedness is preserved; tangent surrogate is
curvature-invariant.}
Figure~\ref{fig:descent}(d) tracks the $\delta$-balancedness defect of
Definition~\ref{def:balanced} along training. Starting from an exactly
balanced initialization, the defect remains below $\sim\!3\times10^{-3}$
throughout at every curvature, so the structural hypothesis of
Theorem~\ref{thm:linear-descent} holds along the trajectory rather than
only at initialization; this is the curvature analogue of the
balancedness experiments of \cite{arora2018convergence}. As a control
(not shown), training the tangent-space surrogate $\mathcal E$ on the
tangent-lifted data yields a per-step contraction exponent that is
identical across all $c$ to machine precision, confirming the reduction of
Section~\ref{sec:disc-tangent}: the surrogate is curvature-blind and
curvature enters only through the dataset.

\begin{figure}[t]
\centering
\includegraphics[width=\linewidth]{figure_collapse.pdf}
\caption{Curvature-collapse check at fixed step size. The intrinsic-loss
convergence exponent $|b|$ (per step) equals the Euclidean tangent-surrogate
exponent exactly at $c=0$ (dashed line) and increases for $c>0$. The
increase is the gradient-amplification effect of $S_c(R)^2$ described in
the text and is the fixed-step counterpart of the shrinking stable step
size in Figure~\ref{fig:descent}(b); the two panels are two views of one
mechanism. The exact coincidence at $c=0$ is the $\mathcal L_c \to
\mathcal E$ collapse of Section~\ref{sec:disc-intrinsic}.}
\label{fig:collapse}
\end{figure}

\paragraph{Scope and limitations.}
These experiments validate the \emph{scaling and qualitative} content of
the descent theory, not its constants: $L_{\mathrm{net}}$ and
$L_{\kappa,R,D}$ are upper bounds, so $\eta^\star$ tracks $1/L_{\mathrm{net}}$
only up to constants, and Figure~\ref{fig:descent}(c) is interpreted in that
light. The setting is deep-linear and bias-free with a single output
manifold (Sections~\ref{sec:disc-hyperbolic}--\ref{sec:disc-intrinsic});
nonlinear activations and M\"obius biases (Section~\ref{sec:future})
are out of scope, as is the layerwise generalization bound of
Theorem~\ref{thm:layerwise-generalization} and its transductive
specialization to Q-GCN
(Corollary~\ref{thm:prgcn_transductive_generalization}), whose empirical
study on a Q-GCN node-classification task we leave to future work.
The reported runs use a single random seed and a fixed radius $R=0.8$
within the regime of the bounded-ball comparison; aggregating $\eta^\star$
over seeds and radii is a straightforward extension.

% =====================================================================
\section{Conclusion}
\label{sec:conclusion}
% =====================================================================

We have organized our results by model class. The \emph{linear case}
(Section~\ref{sec:main-linear}) gives a single-step descent inequality of
Arora--Cohen--Golowich--Hu form for a deep linear network on a Riemannian
output manifold, with a curvature-modulated step-size restriction. The
\emph{nonlinear case} (Section~\ref{sec:main-nonlinear}) extends both
optimization and generalization to the general layerwise Riemannian
neural network model: a curvature-dependent single-step descent lemma
built from recursively defined layerwise smoothness constants, and a
curvature-adaptive generalization bound in which the network's spatial
Lipschitz constant
$\Lambda_N = \prod_j \alpha_j P_j \beta_{j-1}$ is built from the same
layerwise constants. The generalization bound refines
\cite{sarkar2025learning} by exposing the contribution of every
hidden-layer curvature. The \emph{graph case} (Section~\ref{sec:main-graph})
specializes to the pseudo-Riemannian Q-GCN of \cite{xiong2022pseudo},
giving a curvature-adaptive transductive generalization bound as a
corollary. Detailed discussions of how each result specializes to the
tangent-space surrogate, the intrinsic squared hyperbolic distance, and
the pseudo-Riemannian setting are given in place;
Section~\ref{sec:future} collects open questions.

\bibliography{references}
\bibliographystyle{aaai}

% =====================================================================
% Switch to single-column layout for the appendix (proofs render better
% in a single wide column for long display math).
\onecolumn
\appendix
% =====================================================================

% =====================================================================
\section{Proofs}
\label{sec:proofs}
% =====================================================================

\subsection{Proof of Theorem~\ref{thm:linear-descent}}
\label{sec:proof-linear-descent}

We split the proof into the geometric part and the factorization part.

\paragraph{Step 1: The Riemannian pullback gradient.}
For each sample define \(\phi_i(p):=\frac12 d_N(p,y_i)^2\),
\(p\in \mathcal M_N\). On the chosen geodesically convex tube,
\(\nabla_p \phi_i(p) = -\log_p(y_i)\). Let \(v_i=A\xi_i\),
\(p_i=\exp_{o_N}(v_i)\). Then
\(\mathcal F(A) = \frac1m\sum_{i=1}^m \phi_i(\exp_{o_N}(A\xi_i))\).
By the chain rule,
\[
\nabla_A\mathcal F(A)
=
\frac1m\sum_{i=1}^m
\left[D\exp_{o_N}(v_i)\right]^* \nabla_p\phi_i(p_i) \,\xi_i^*.
\]
Since \(\nabla_p\phi_i(p_i)=-\log_{p_i}(y_i)\),
\begin{equation}
\label{eq:pullback-gradient}
G
=
\nabla_A\mathcal F(A)
=
-\frac1m\sum_{i=1}^m
\left[D\exp_{o_N}(A\xi_i)\right]^* \log_{p_i(A)}(y_i) \,\xi_i^*.
\end{equation}
In Euclidean space, \(D\exp_{o_N}=I\) and \(\log_p(y)=y-p\), so this
reduces to the usual residual gradient.

\paragraph{Step 2: Curvature-dependent smoothness of the pullback loss.}
Let \(\Delta:V_0\to V_N\) be a perturbation of \(A\), and define
\(u_i:=\Delta \xi_i\). Differentiating twice,
\[
\begin{aligned}
D^2\mathcal F(A)[\Delta,\Delta]
=
\frac1m\sum_{i=1}^m
\Big[
&\nabla_p^2\phi_i(p_i)
\big(D\exp_{o_N}(v_i)u_i, D\exp_{o_N}(v_i)u_i\big)
\\
&+ \left\langle \nabla_p\phi_i(p_i),
D^2\exp_{o_N}(v_i)[u_i,u_i] \right\rangle_{p_i}
\Big].
\end{aligned}
\]
On the tube, \(\|D\exp_{o_N}(v_i)\|_{\mathrm{op}} \leq S_\kappa(R)\),
\(\|\nabla_p^2\phi_i(p_i)\|_{\mathrm{op}} \leq H_\kappa(D)\),
\(\|\nabla_p\phi_i(p_i)\| = d_N(p_i,y_i) \leq D\),
and \(\|D^2\exp_{o_N}(v_i)\|_{\mathrm{op}} \leq B_\kappa(R)\).
Therefore
\[
|D^2\mathcal F(A)[\Delta,\Delta]|
\leq
\frac1m\sum_{i=1}^m \|\xi_i\|^2
\left[S_\kappa(R)^2H_\kappa(D) + D\,B_\kappa(R)\right]
\|\Delta\|_{\mathrm{HS}}^2
=
L_{\kappa,R,D} \|\Delta\|_{\mathrm{HS}}^2.
\]
Hence \(\mathcal F\) is \(L_{\kappa,R,D}\)-smooth on the tube:
\[
\mathcal F(B)
\leq
\mathcal F(A)
+
\langle G,B-A\rangle
+
\frac{L_{\kappa,R,D}}2 \|B-A\|_{\mathrm{HS}}^2
\]
for all \(A,B\) whose connecting segment remains in \(\mathcal T_{R,D}\).

\paragraph{Step 3: Gradients with respect to the factors.}
Since \(A=W_N\cdots W_1\), a variation \(\Delta W_j\) produces
\(\Delta A = W_{j+1:N}\Delta W_jW_{1:j-1}\). Therefore
\[
D_{W_j}\mathcal L[\Delta W_j]
=
\left\langle G, W_{j+1:N}\Delta W_jW_{1:j-1} \right\rangle_{\mathrm{HS}}
=
\left\langle W_{j+1:N}^*GW_{1:j-1}^*, \Delta W_j \right\rangle_{\mathrm{HS}}.
\]
Thus \(\nabla_{W_j}\mathcal L = W_{j+1:N}^*GW_{1:j-1}^*.\)

\paragraph{Step 4: Expansion of the end-to-end update.}
The gradient step is \(W_j^+ = W_j - \eta W_{j+1:N}^*GW_{1:j-1}^*\).
Expanding \(A^+=W_N^+\cdots W_1^+\) to first order in \(\eta\),
\[
A^+-A = -\eta \mathcal P_A(G) + R_\eta,
\quad
\mathcal P_A(G)
:=
\sum_{j=1}^N
W_{j+1:N}W_{j+1:N}^* G W_{1:j-1}^*W_{1:j-1}.
\]
The term \(R_\eta\) contains only second-order and higher-order products
of the perturbations \(\Delta W_j:=W_j^+-W_j\). Using
\(\|W_j\|_{\mathrm{op}}\leq M\), there exists \(C_N>0\), depending only on
\(N\), such that
\[
\|A^+-A\|_{\mathrm{HS}} \leq C_N\eta M^{2N-2}\|G\|_{\mathrm{HS}},
\qquad
|\langle G,R_\eta\rangle|
\leq
C_N\eta^2M^{3N-4}G_{\kappa,R,D} \|G\|_{\mathrm{HS}}^2.
\]

\paragraph{Step 5: The first-order descent term.}
We compute
\[
\langle G,\mathcal P_A(G)\rangle_{\mathrm{HS}}
=
\sum_{j=1}^N \|W_{j+1:N}^*GW_{1:j-1}^*\|_{\mathrm{HS}}^2.
\]
By the assumed balancedness/coercivity condition,
\(\langle G,\mathcal P_A(G)\rangle_{\mathrm{HS}} \geq \tfrac34
\sigma_{\min}(A)^{2(N-1)/N} \|G\|_{\mathrm{HS}}^2\). Therefore
\[
\langle G,A^+-A\rangle_{\mathrm{HS}}
\leq
- \tfrac34\eta \sigma_{\min}(A)^{2(N-1)/N} \|G\|_{\mathrm{HS}}^2
+
C_N\eta^2M^{3N-4}G_{\kappa,R,D} \|G\|_{\mathrm{HS}}^2.
\]

\paragraph{Step 6: Apply curvature-dependent smoothness.}
By \(L_{\kappa,R,D}\)-smoothness,
\[
\mathcal F(A^+)
\leq
\mathcal F(A)
+
\langle G,A^+-A\rangle_{\mathrm{HS}}
+
\frac{L_{\kappa,R,D}}2 \|A^+-A\|_{\mathrm{HS}}^2.
\]
Using the bounds above,
\[
\begin{aligned}
\mathcal F(A^+)-\mathcal F(A)
&\leq
- \tfrac34\eta \sigma_{\min}(A)^{2(N-1)/N} \|G\|_{\mathrm{HS}}^2
\\
&\qquad
+ C_N\eta^2M^{3N-4}G_{\kappa,R,D} \|G\|_{\mathrm{HS}}^2
+ C_N\eta^2L_{\kappa,R,D}M^{4N-4} \|G\|_{\mathrm{HS}}^2.
\end{aligned}
\]
The step-size assumption gives
\[
C_N\eta \left( M^{3N-4}G_{\kappa,R,D} + L_{\kappa,R,D}M^{4N-4} \right)
\leq \tfrac14 \sigma_{\min}(A)^{2(N-1)/N}.
\]
Hence the two positive second-order terms are at most \(\tfrac14\eta
\sigma_{\min}(A)^{2(N-1)/N} \|G\|_{\mathrm{HS}}^2\). Therefore
\[
\mathcal F(A^+)-\mathcal F(A)
\leq
- \tfrac12\eta \sigma_{\min}(A)^{2(N-1)/N} \|G\|_{\mathrm{HS}}^2.
\]
Since \(G=\nabla_A\mathcal F(A)\), this is the stated bound. Finally,
because \(\mathcal L(W_1,\dots,W_N)=\mathcal F(W_N\cdots W_1)\), the same
inequality reads
\[
\mathcal L(W_1^+,\dots,W_N^+) - \mathcal L(W_1,\dots,W_N)
\leq
- \tfrac{\eta}{2} \sigma_{\min}(A)^{2(N-1)/N}
\|\nabla_A\mathcal F(A)\|_{\mathrm{HS}}^2.
\qquad\square
\]

\subsection{Proof of Theorem~\ref{thm:general-descent}}
\label{sec:proof-general-descent}

Fix a sample \(i\) and define \(f_i(\theta) := \ell_i(h_N^{(i)}(\theta))\).
By the chain rule,
\[
D_\theta f_i
=
D\ell_i(h_N^{(i)})
\circ
D_\theta h_N^{(i)}.
\]
Differentiating again,
\[
\begin{aligned}
D_\theta^2 f_i[\Delta,\Delta]
&=
\nabla_p^2\ell_i(h_N^{(i)})
\left[
D_\theta h_N^{(i)}[\Delta],
D_\theta h_N^{(i)}[\Delta]
\right]
\\
&\qquad
+
\left\langle
\nabla_p\ell_i(h_N^{(i)}),
D_\theta^2 h_N^{(i)}[\Delta,\Delta]
\right\rangle_{g_N}.
\end{aligned}
\]
Using the definitions of \(G_N,H_N,J_N,K_N\), we obtain
\[
\begin{aligned}
|D_\theta^2 f_i[\Delta,\Delta]|
&\le
H_N \|D_\theta h_N^{(i)}[\Delta]\|^2
+
G_N \|D_\theta^2 h_N^{(i)}[\Delta,\Delta]\|
\\
&\le
H_NJ_N^2\|\Delta\|^2 + G_NK_N\|\Delta\|^2
\\
&=
L_{\mathrm{net}}\|\Delta\|^2.
\end{aligned}
\]
Averaging over \(i=1,\dots,m\), the empirical loss
\(\mathcal L(\theta) = \frac1m\sum_{i=1}^m f_i(\theta)\) is
\(L_{\mathrm{net}}\)-smooth on the tube. Therefore
\[
\mathcal L(\theta^+)
\le
\mathcal L(\theta)
+
\left\langle
\nabla_\theta\mathcal L(\theta), \theta^+-\theta
\right\rangle
+
\frac{L_{\mathrm{net}}}{2}
\|\theta^+-\theta\|^2.
\]
Since \(\theta^+-\theta = -\eta\nabla_\theta\mathcal L(\theta)\),
\[
\mathcal L(\theta^+)
\le
\mathcal L(\theta)
-
\eta \left(1-\frac{\eta L_{\mathrm{net}}}{2}\right)
\|\nabla_\theta\mathcal L(\theta)\|^2.
\]
If \(\eta\le 1/L_{\mathrm{net}}\) (i.e.~\eqref{eq:general-stepsize}), then
\(1-\eta L_{\mathrm{net}}/2 \ge 1/2\), so
\[
\mathcal L(\theta^+)
\le
\mathcal L(\theta) - \frac{\eta}{2} \|\nabla_\theta\mathcal L(\theta)\|^2,
\]
which is~\eqref{eq:general-descent}. $\square$

\subsection{Proof of Theorem~\ref{thm:layerwise-generalization}}
\label{sec:proof-layerwise-gen}

The argument is a layerwise refinement of \cite{sarkar2025learning}:
Steps~1--2 collapse the network to a single $L_\ell \Lambda_N$-Lipschitz
function on $\mathcal M_0$ using the recursive constants of
Section~\ref{sec:prelim-recursive}; Steps~3--5 then apply the
Bishop--Gromov covering, function-class covering, and Dudley entropy
integral of \cite[Theorems~3.1--3.3]{sarkar2025learning}; Step~6 is the
standard symmetrization and bounded-difference concentration.

\paragraph{Step 1: Layerwise spatial Lipschitz of the network.}
For each $j$, we claim
\(
\sup_{x\in\mathcal T_0,\theta}\|D_x h_j(x;\theta)\|_{\mathrm{op}} \le \Lambda_j,
\)
where $\Lambda_j$ is given by~\eqref{eq:Lambda-recursive}.
The case $j=0$ is trivial. Inductively, the manifold
layer~\eqref{eq:phi-layer} gives, by the chain rule,
\[
D_x h_j(x;\theta)
=
D\exp_{o_j}^{(j)}\!\cdot\!D_u F_j \!\cdot\!D\log_{o_{j-1}}^{(j-1)}\!\cdot\! D_x h_{j-1}(x;\theta).
\]
Taking operator norms and using
$\|D\exp_{o_j}^{(j)}\|_{\mathrm{op}}\le \alpha_j$,
$\|D_u F_j\|_{\mathrm{op}}\le P_j$,
$\|D\log_{o_{j-1}}^{(j-1)}\|_{\mathrm{op}}\le \beta_{j-1}$,
the inductive bound follows. Setting $j=N$ gives
$\|D_x h_N(x;\theta)\|_{\mathrm{op}}\le \Lambda_N$, hence
$h_N(\cdot;\theta) : \mathcal M_0 \to \mathcal M_N$ is
$\Lambda_N$-Lipschitz with respect to the Riemannian metrics; the passage
from operator norms on tangent spaces to geodesic distances on the
compact tubes $\mathcal T_0, \mathcal T_N$ costs only a bounded factor,
absorbed into $\Lambda_N$.

\paragraph{Step 2: The loss class on $\mathcal M_0$.}
Define $f_\theta(x):=\ell(h_N(x;\theta))$ and
$\mathcal F := \{f_\theta : \theta \text{ in the parameter tube}\}$. By
hypothesis $\ell$ is $L_\ell$-Lipschitz on $\mathcal T_N$ and bounded by
$B$, so by Step~1 each $f_\theta$ is $L_\ell \Lambda_N$-Lipschitz on
$\mathcal M_0$ and uniformly bounded by $B$. Thus $\mathcal F$ embeds in
the class of $L_\ell \Lambda_N$-Lipschitz $\R$-valued functions on
$\mathcal M_0$ with $\|f\|_\infty \le B$.

\paragraph{Step 3: Curvature-adaptive covering of $\mathcal M_0$.}
By Bishop--Gromov volume comparison
\cite[Theorem~3.1]{sarkar2025learning}, for any
$\varepsilon < \tfrac12 \iota_0$,
\[
\log \mathcal N(\mathcal M_0, \varepsilon, d_{\mathcal M_0})
\;\le\;
\log\!\left(\frac{\mathrm{Vol}(\mathcal M_0)}{\inf_{x}\mathrm{Vol}(B(x,\varepsilon/2))}\right)
+ C\sqrt{|\kappa_{\max}^{(0)}|}\,\varepsilon.
\]
For $\kappa_{\max}^{(0)} \le 0$, Bishop--Gromov gives
$\inf_x \mathrm{Vol}(B(x,\varepsilon/2)) \ge
\omega_{d_0}(\varepsilon/2)^{d_0} e^{-c\sqrt{|\kappa_{\max}^{(0)}|}\,\varepsilon}$,
and the bound becomes
\(
\log\mathcal N(\mathcal M_0,\varepsilon,d_{\mathcal M_0})
\le d_0 \log(c'/\varepsilon)
+ C\sqrt{|\kappa_{\max}^{(0)}|}\,\varepsilon
+ \log(\mathrm{Vol}(\mathcal M_0)/\omega_{d_0}).
\)

\paragraph{Step 4: Covering of the loss class.}
By the function-class covering bound of
\cite[Theorem~3.2]{sarkar2025learning}, for
$\varepsilon < \iota_0/(2L_\ell\Lambda_N)$,
\[
\log \mathcal N(\mathcal F, \varepsilon, \|\cdot\|_\infty)
\;\le\;
\mathcal N\!\left(\mathcal M_0,\tfrac{\varepsilon}{2L_\ell\Lambda_N},d_{\mathcal M_0}\right)
\cdot \log\!\frac{4B}{\varepsilon}.
\]
Substituting Step~3,
\[
\log \mathcal N(\mathcal F, \varepsilon, \|\cdot\|_\infty)
\;\lesssim\;
\Big[
d_0 \log\!\tfrac{2L_\ell\Lambda_N}{\varepsilon}
+ \sqrt{|\kappa_{\max}^{(0)}|}\,\tfrac{\varepsilon}{L_\ell\Lambda_N}
+ \log\!\tfrac{\mathrm{Vol}(\mathcal M_0)}{\omega_{d_0}}
\Big]\log\!\frac{4B}{\varepsilon}.
\]

\paragraph{Step 5: Dudley entropy integral and Rademacher complexity.}
By Dudley's entropy bound, the empirical Rademacher complexity of
$\mathcal F$ satisfies
\[
\mathrm{Rad}_n(\mathcal F)
\;\le\;
\inf_{\alpha>0}
\Big\{4\alpha + \tfrac{12}{\sqrt n}\int_\alpha^B
\sqrt{\log\mathcal N(\mathcal F,\varepsilon,\|\cdot\|_\infty)}\,d\varepsilon\Big\}.
\]
Substituting Step~4 and choosing $\alpha = n^{-1/d_0}$ yields (the
computation is identical to \cite[Theorem~3.3, Steps~3--5]{sarkar2025learning},
with $L$ replaced by $L_\ell\Lambda_N$)
\[
\mathrm{Rad}_n(\mathcal F)
\;\le\;
O\!\left(
\frac{\sqrt{d_0 \log(L_\ell\Lambda_N\sqrt n)}+\psi(\kappa_{\max}^{(0)}, L_\ell\Lambda_N)}{n^{1/d_0}}
\right),
\]
with $\psi(\kappa,L)=\sqrt{|\kappa|}/L$ for $\kappa<0$ and $\psi=0$ for
$\kappa\ge 0$.

\paragraph{Step 6: Concentration.}
By symmetrization and McDiarmid's inequality applied to the bounded loss
class (see, e.g., \cite[Corollary~3.4]{sarkar2025learning}), with
probability at least $1-\delta$, for all $\theta$ in the parameter tube,
\[
|R(\theta)-\hat R_n(\theta)|
\le 2L_\ell\,\mathrm{Rad}_n(\mathcal F)
+ 3B\sqrt{\tfrac{\log(2/\delta)}{2n}}.
\]
Combining with Step~5 and absorbing constants gives the stated bound.
$\square$

\subsection{Proof of Corollary~\ref{thm:prgcn_transductive_generalization}}
\label{sec:proof-gen}

The proof applies Theorem~\ref{thm:layerwise-generalization} after two
adjustments specific to the Q-GCN setting: the diffeomorphic chart
$\psi$ to the geodesically complete product manifold $M_\psi$, and the
transductive Rademacher comparison of \cite{tang2023towards} in place of
the standard PAC concentration step.

\paragraph{Step 1: Transductive reduction.}
By the transductive Rademacher comparison of
\cite[Theorem~4.3]{tang2023towards}, with probability at least
$1 - \delta/2$,
\begin{equation}
\label{eq:pr_step1}
R_u(w^{(T+1)}) \leq R_m(w^{(T+1)}) + \mathfrak R_{m+u}(\mathcal F_Q)
+ c_0 Q \sqrt{\min(m, u)} + \sqrt{\tfrac{SQ}{2} \log\tfrac{2}{\delta}},
\end{equation}
where $Q = \tfrac1m + \tfrac1u$ and $S$ is the same transductive
combinatorial factor. It remains to upper bound the transductive
Rademacher complexity $\mathfrak R_{m+u}(\mathcal F_Q)$.

\paragraph{Step 2: Dudley-type entropy integral.}
Exactly as in \cite{tang2023towards}, the Rademacher complexity is bounded
by
\begin{equation}
\label{eq:pr_step2}
\mathfrak R_{m+u}(\mathcal F_Q) \leq \tfrac{(m+u)^{3/2}}{mu}
\left[ b_\ell + 12 \int_0^{M_Q}
\sqrt{\log \mathcal N(r, \mathcal F_Q, d_{HS})}\, dr \right],
\end{equation}
where $M_Q$ is the diameter of the loss class around initialization under
the empirical metric $d_{HS}$. The only step that changes is the
covering-number estimate.

\paragraph{Step 3: Covering by the diffeomorphic product manifold.}
The Q-GCN forward map is a composition of (a) the chart $\psi$ from
$\mathcal Q_\beta^{s,t}$ to $\mathbb S^t_{-\beta} \times \R^s$;
(b) Lipschitz tangent-space linear maps;
(c) neighborhood summation;
(d) a pointwise nonlinearity $\sigma$;
(e) the inverse chart $\psi^{-1}$ and broken-case bias translation.
Because $\psi, \psi^{-1}$ are smooth diffeomorphisms and $o$ is chosen so
that the tangent space of the diffeomorphic model matches the tangent
space of the pseudo-hyperboloid itself, the Q-GCN loss class is controlled
by covering the product manifold $M_\psi = \prod_\ell (\mathbb
S^t_{-\beta_\ell} \times \R^{s_\ell})$ and discretizing function values on
$M_\psi$.

\paragraph{Step 4: Curvature-adaptive covering of $M_\psi$.}
For each spherical factor $\mathbb S^t_{-\beta_\ell}$, by the standard
curvature-adaptive manifold covering bound
\cite[e.g.][]{niyogi2008finding},
\[
\log \mathcal N(\mathbb S^t_{-\beta_\ell}, \varepsilon, d_g)
\;\leq\; \log\!\left(\tfrac{\mathrm{Vol}(\mathbb S^t_{-\beta_\ell})}
{\inf_x \mathrm{Vol}(B(x, \varepsilon/2))}\right)
+ C_\ell \sqrt{|\kappa_{\max}^{(\ell)}|}\, \varepsilon.
\]
The Euclidean factor $\R^{s_\ell}$ contributes standard polynomial entropy
on the bounded region traversed by the iterates. Combining,
\begin{equation}
\label{eq:pr_step4}
\log \mathcal N(M_\psi, \varepsilon, d_g)
\;\leq\; d_{\mathrm{eff}} \log\tfrac{c}{\varepsilon} + C \Psi_{\mathrm{curv}}
+ \log\!\big(\mathrm{Vol}(M_\psi)/\omega_{d_{\mathrm{eff}}}\big),
\end{equation}
for a constant $c > 0$ depending on the radius of the forward trajectory
and the layerwise diffeomorphic charts, with $\Psi_{\mathrm{curv}} =
\sum_\ell c_\ell \sqrt{|\kappa_{\max}^{(\ell)}|}$.

\paragraph{Step 5: Covering the Q-GCN loss class.}
Discretize function values on an $\varepsilon/(2L_Q)$-net of $M_\psi$:
\begin{equation}
\label{eq:pr_step5}
\log \mathcal N(\mathcal F_Q, \varepsilon, \|\cdot\|_\infty)
\;\leq\;
\mathcal N(M_\psi, \tfrac{\varepsilon}{2L_Q}, d_g) \log\tfrac{4B}{\varepsilon}.
\end{equation}
Combining \eqref{eq:pr_step4} and \eqref{eq:pr_step5},
\begin{equation}
\label{eq:pr_step6}
\log \mathcal N(\mathcal F_Q, \varepsilon, \|\cdot\|_\infty)
\;\lesssim\;
\Big[ d_{\mathrm{eff}} \log\tfrac{2L_Q}{\varepsilon} + \Psi_{\mathrm{curv}}
+ \log(\mathrm{Vol}(M_\psi)/\omega_{d_{\mathrm{eff}}}) \Big]
\log\tfrac{4B}{\varepsilon}.
\end{equation}
Since $d_{HS} \leq \|\cdot\|_\infty$, the same bound applies to
$\log \mathcal N(r, \mathcal F_Q, d_{HS})$.

\paragraph{Step 6: Entropy integral estimate.}
Substituting \eqref{eq:pr_step6} into \eqref{eq:pr_step2} and integrating
as in the original proof yields
\[
\mathfrak R_{m+u}(\mathcal F_Q) \leq
C' \tfrac{(m+u)^{3/2}}{mu} L_Q
\big[ d_{\mathrm{eff}} \log(L_Q (m+u)^{1/d_{\mathrm{eff}}}) + \Psi_{\mathrm{curv}}
+ \log(\mathrm{Vol}(M_\psi)/\omega_{d_{\mathrm{eff}}}) \big]^{1/2},
\]
which is $\mathfrak C_Q$ from the theorem statement.

\paragraph{Step 7: Optimization radius and conclusion.}
The stochastic-optimization radius bound on $\|w^{(T+1)} - w^{(1)}\|_2$
inherited from \cite[Theorem~4.3]{tang2023towards} gives the three regimes
of $T^*_\alpha(T)$. Inserting the Rademacher bound into~\eqref{eq:pr_step1}
completes the proof. $\square$

\end{document}